\input{preamble}

% Bon, commençons ici.
%
\begin{document}

\title{Propriétés des champs algébriques}


\maketitle

\phantomsection
\label{section-phantom}

\tableofcontents

\section{Introduction}
\label{section-introduction}

\noindent
Voir
Algebraic Stacks, Section \ref{algebraic-section-introduction},
pour une brève introduction aux champs algébriques, et lire
une partie de ce chapitre pour connaître les fondements que nous adoptons.
Dans ce chapitre-là, nous nous attachons à distinguer soigneusement
les schémas, les espaces algébriques et les champs algébriques ; à partir
du présent chapitre, nous emploierons l'abus de langage usuel selon lequel
toutes ces notions s'emploient indifféremment.

\medskip\noindent
L'objectif de ce chapitre est d'introduire quelques notions et
propriétés fondamentales des champs algébriques. Une référence
essentielle pour le cas des champs algébriques quasi-séparés à diagonale
représentable est \cite{LM-B}.



\section{Conventions et abus de langage}
\label{section-conventions}

\noindent
Fixons un grand site fppf $\Sch_{fppf}$.
Tous les schémas appartiennent à $\Sch_{fppf}$.
Et tous les anneaux $A$ considérés ont la propriété que
$\Spec(A)$ est (isomorphe) à un objet de ce grand site.

\medskip\noindent
Fixons également un schéma de base $S$, qui, d'après les conventions ci-dessus, est un élément
de $\Sch_{fppf}$. Le lecteur qui ne s'intéresse qu'au
cas absolu peut prendre $S = \Spec(\mathbf{Z})$.

\medskip\noindent
Voici nos conventions concernant les champs algébriques :
\begin{enumerate}
\item Lorsque nous dirons {\it champ algébrique}, nous entendrons un champ
algébrique au-dessus de $S$, c'est-à-dire une catégorie fibrée en groupoïdes
$p : \mathcal{X} \to (\Sch/S)_{fppf}$
qui satisfait aux conditions de
Algebraic Stacks, Définition \ref{algebraic-definition-algebraic-stack}.
\item Nous dirons que $f : \mathcal{X} \to \mathcal{Y}$ est un
{\it morphisme de champs algébriques} pour désigner un $1$-morphisme
de champs algébriques au-dessus de $S$, c'est-à-dire un $1$-morphisme de catégories fibrées
en groupoïdes au-dessus de $(\Sch/S)_{fppf}$, voir
Algebraic Stacks,
Définition \ref{algebraic-definition-morphism-algebraic-stacks}.
\item Un {\it $2$-morphisme} $\alpha : f \to g$ désignera
un $2$-morphisme dans la $2$-catégorie des champs algébriques au-dessus de
$S$, voir
Algebraic Stacks,
Définition \ref{algebraic-definition-morphism-algebraic-stacks}.
\item Étant donnés des morphismes $\mathcal{X} \to \mathcal{Z}$
et $\mathcal{Y} \to \mathcal{Z}$ de champs algébriques,
nous appelons abusivement le $2$-produit fibré
$\mathcal{X} \times_\mathcal{Z} \mathcal{Y}$ le {\it produit fibré}.
\item Nous noterons $\mathcal{X} \times_S \mathcal{Y}$ le
produit des champs algébriques $\mathcal{X}$, $\mathcal{Y}$.
\item Nous abuserons souvent de la notation et dirons que deux champs algébriques
$\mathcal{X}$ et $\mathcal{Y}$ sont {\it isomorphes} s'ils sont
équivalents dans cette $2$-catégorie.
\end{enumerate}
Voici nos conventions concernant les espaces algébriques.
\begin{enumerate}
\item Si nous disons que $X$ est un {\it espace algébrique}, nous entendons que
$X$ est un espace algébrique au-dessus de $S$, c'est-à-dire que $X$ est un préfaisceau sur
$(\Sch/S)_{fppf}$ qui satisfait aux conditions de
Spaces, Définition \ref{spaces-definition-algebraic-space}.
\item Un {\it morphisme d'espaces algébriques} $f :X \to Y$ est un morphisme
d'espaces algébriques au-dessus de $S$ tel que défini dans
Spaces, Définition \ref{spaces-definition-morphism-algebraic-spaces}.
\item Nous ne distinguerons {\bf pas} un espace algébrique $X$
du champ algébrique $\mathcal{S}_X \to (\Sch/S)_{fppf}$
auquel il donne lieu, voir
Algebraic Stacks, Lemme \ref{algebraic-lemma-representable-algebraic}.
\item En particulier, un {\it morphisme} $f : X \to \mathcal{Y}$ de $X$
vers un champ algébrique $\mathcal{Y}$ désigne un morphisme
$f : \mathcal{S}_X \to \mathcal{Y}$ de champs algébriques.
Il en va de même pour les morphismes $\mathcal{Y} \to X$.
\item De plus, étant donné un champ algébrique $\mathcal{X}$, nous dirons que
{\it $\mathcal{X}$ est un espace algébrique} pour indiquer que $\mathcal{X}$
est représentable par un espace algébrique, voir
Algebraic Stacks,
Définition \ref{algebraic-definition-representable-by-algebraic-space}.
\item Nous emploierons la convention de notation suivante : si nous
désignons un champ algébrique par une majuscule romaine
(telle que $X, Y, Z, A, B, \ldots$), son
champ d'inertie sera trivial, et ce sera donc un espace algébrique, voir
Algebraic Stacks,
Proposition \ref{algebraic-proposition-algebraic-stack-no-automorphisms}.
\end{enumerate}
Voici nos conventions concernant les schémas.
\begin{enumerate}
\item Si nous disons que $X$ est un {\it schéma}, nous entendons que
$X$ est un schéma au-dessus de $S$, c'est-à-dire que $X$ est un objet de $(\Sch/S)_{fppf}$.
\item Par {\it morphisme de schémas}, nous entendons un morphisme de schémas au-dessus de $S$.
\item Nous ne distinguerons {\bf pas} un schéma $X$ du
champ algébrique $\mathcal{S}_X \to (\Sch/S)_{fppf}$ auquel il donne lieu,
voir
Algebraic Stacks, Lemme \ref{algebraic-lemma-representable-algebraic}.
\item En particulier, un {\it morphisme} $f : X \to \mathcal{Y}$ d'un
schéma $X$ vers un champ algébrique $\mathcal{Y}$ désigne un morphisme
$f : \mathcal{S}_X \to \mathcal{Y}$ de champs algébriques.
Il en va de même pour les morphismes $\mathcal{Y} \to X$.
\item De plus, étant donné un champ algébrique $\mathcal{X}$, nous dirons que
{\it $\mathcal{X}$ est un schéma} pour indiquer que $\mathcal{X}$
est représentable, voir
Algebraic Stacks, Section \ref{algebraic-section-representable}.
\end{enumerate}
Voici nos conventions concernant les morphismes de champs algébriques :
\begin{enumerate}
\item Un morphisme $f : \mathcal{X} \to \mathcal{Y}$
de champs algébriques est {\it représentable}, ou
{\it représentable par des schémas}, si, pour tout schéma
$T$ et tout morphisme $T \to \mathcal{Y}$, le produit fibré
$T \times_\mathcal{Y} \mathcal{X}$ est un schéma.
Voir
Algebraic Stacks, Section \ref{algebraic-section-representable-morphism}.
\item Un morphisme $f : \mathcal{X} \to \mathcal{Y}$
de champs algébriques est {\it représentable par des espaces algébriques}
si, pour tout schéma $T$ et tout morphisme $T \to \mathcal{Y}$, le produit fibré
$T \times_\mathcal{Y} \mathcal{X}$ est un espace algébrique.
Voir Algebraic Stacks,
Définition \ref{algebraic-definition-representable-by-algebraic-spaces}.
Dans ce cas, $Z \times_\mathcal{Y} \mathcal{X}$ est un espace
algébrique dès que $Z \to \mathcal{Y}$ est un morphisme dont la source est
un espace algébrique, voir
Algebraic Stacks,
Lemme \ref{algebraic-lemma-base-change-by-space-representable-by-space}.
\item Nous pourrons abuser du langage et dire qu'un diagramme de champs algébriques
commute si le diagramme est $2$-commutatif dans la $2$-catégorie des
champs algébriques.
\end{enumerate}
Remarquons que tout morphisme $X \to \mathcal{Y}$ d'un espace algébrique
vers un champ algébrique est représentable par des espaces algébriques, voir
Algebraic Stacks, Lemme \ref{algebraic-lemma-representable-diagonal}.
Nous utiliserons ce résultat fondamental sans plus le rappeler.




\section{Propriétés des morphismes représentables par des espaces algébriques}
\label{section-properties-morphisms}

\noindent
Nous étudierons les propriétés des morphismes (quelconques) de champs algébriques
dans un chapitre qui leur est propre. Pour les morphismes représentables par des espaces algébriques, nous savons ce que
signifie être surjectif, lisse, ou \'etale, etc. Cela s'applique en particulier
aux morphismes $X \to \mathcal{Y}$ des espaces algébriques vers les champs algébriques.
Dans cette section, nous rappelons le fonctionnement de cette notion, nous dressons la liste des propriétés
auxquelles elle s'applique et nous démontrons quelques lemmes faciles.

\medskip\noindent
Notre premier lemme affirme qu'un morphisme est représentable par des espaces algébriques
s'il le devient après un changement de base par un morphisme plat,
localement de présentation finie et surjectif.

\begin{lemma}
\label{lemma-check-representable-covering}
Soit $f : \mathcal{X} \to \mathcal{Y}$ un morphisme de champs algébriques.
Soit $W$ un espace algébrique et soit $W \to \mathcal{Y}$ un morphisme surjectif,
localement de présentation finie et plat. Les assertions suivantes sont équivalentes :
\begin{enumerate}
\item $f$ est représentable par des espaces algébriques, et
\item $W \times_\mathcal{Y} \mathcal{X}$ est un espace algébrique.
\end{enumerate}
\end{lemma}

\begin{proof}
L'implication (1) $\Rightarrow$ (2) est
Algebraic Stacks,
Lemme \ref{algebraic-lemma-base-change-by-space-representable-by-space}.
Réciproquement, soit $W \to \mathcal{Y}$ comme en (2). Pour démontrer (1), il
suffit de montrer que $f$ est fidèle sur les catégories fibres, voir
Algebraic Stacks,
Lemme \ref{algebraic-lemma-characterize-representable-by-algebraic-spaces}.
L'hypothèse (2) implique en particulier que
$W \times_\mathcal{Y} \mathcal{X} \to W$ est fidèle.
La fidélité de $f$ résulte donc de
Stacks, Lemme \ref{stacks-lemma-faithful-descent}.
\end{proof}

\noindent
Soit $P$ une propriété des morphismes d'espaces algébriques qui est
locale pour la topologie fppf sur le but et préservée par tout changement de base.
Soit $f : \mathcal{X} \to \mathcal{Y}$ un morphisme de champs algébriques
représentable par des espaces algébriques. Nous disons alors que
{\it $f$ possède la propriété $P$} si et seulement si, pour tout schéma $T$
et tout morphisme $T \to \mathcal{Y}$, le morphisme d'espaces algébriques
$T \times_\mathcal{Y} \mathcal{X} \to T$ possède la propriété $P$, voir
Algebraic Stacks,
Définition \ref{algebraic-definition-relative-representable-property}.

\medskip\noindent
Il se trouve que si $f : \mathcal{X} \to \mathcal{Y}$ est représentable
par des espaces algébriques et possède la propriété $P$, alors, pour tout morphisme de champs algébriques
$\mathcal{Y}' \to \mathcal{Y}$, le changement de base
$\mathcal{Y}' \times_\mathcal{Y} \mathcal{X} \to \mathcal{Y}'$
possède la propriété $P$, voir
Algebraic Stacks,
Lemmes \ref{algebraic-lemma-base-change-representable-by-spaces} et
\ref{algebraic-lemma-base-change-representable-transformations-property}.
Si la propriété $P$ est préservée par composition, cela vaut
également pour les morphismes de champs algébriques représentables par
des espaces algébriques, voir
Algebraic Stacks,
Lemmes \ref{algebraic-lemma-composition-representable-by-spaces} et
\ref{algebraic-lemma-composition-representable-transformations-property}.
En outre, dans ce cas, les produits
$\mathcal{X}_1 \times \mathcal{X}_2 \to \mathcal{Y}_1 \times \mathcal{Y}_2$
de morphismes représentables par des espaces algébriques et possédant la propriété $\mathcal{P}$
possèdent la propriété $\mathcal{P}$, voir
Algebraic Stacks, Lemme
\ref{algebraic-lemma-product-representable-transformations-property}.

\medskip\noindent
Enfin, si nous avons deux propriétés $P, P'$ de morphismes d'espaces algébriques
qui sont locales pour la topologie fppf sur le but et préservées par tout changement de base,
et si $P(f) \Rightarrow P'(f)$ pour tout morphisme $f$, alors la même
implication vaut pour les propriétés correspondantes des morphismes de champs algébriques
représentables par des espaces algébriques, voir
Algebraic Stacks, Lemme
\ref{algebraic-lemma-representable-transformations-property-implication}.
Nous l'utiliserons {\it sans autre mention} ci-dessous et dans les
chapitres suivants.

\medskip\noindent
La discussion ci-dessus s'applique à chacune des propriétés suivantes des
morphismes d'espaces algébriques :
\begin{enumerate}
\item quasi-compact, voir
Morphisms of Spaces,
Lemme \ref{spaces-morphisms-lemma-base-change-quasi-compact}
et
Descent on Spaces,
Lemme \ref{spaces-descent-lemma-descending-property-quasi-compact},
\item quasi-séparé, voir
Morphisms of Spaces,
Lemme \ref{spaces-morphisms-lemma-base-change-separated}
et
Descent on Spaces,
Lemme \ref{spaces-descent-lemma-descending-property-quasi-separated},
\item universellement fermé, voir
Morphisms of Spaces,
Lemme \ref{spaces-morphisms-lemma-base-change-universally-closed}
et
Descent on Spaces,
Lemme \ref{spaces-descent-lemma-descending-property-universally-closed},
\item universellement ouvert, voir
Morphisms of Spaces,
Lemme \ref{spaces-morphisms-lemma-base-change-universally-open}
et
Descent on Spaces,
Lemme \ref{spaces-descent-lemma-descending-property-universally-open},
\item universellement submersif, voir
Morphisms of Spaces,
Lemme \ref{spaces-morphisms-lemma-base-change-universally-submersive}
et
Descent on Spaces,
Lemme \ref{spaces-descent-lemma-descending-property-universally-submersive},
\item un homéomorphisme universel, voir
Morphisms of Spaces,
Lemme \ref{spaces-morphisms-lemma-base-change-universal-homeomorphism}
et
Descent on Spaces,
Lemme \ref{spaces-descent-lemma-descending-property-universal-homeomorphism},
\item surjectif, voir
Morphisms of Spaces,
Lemme \ref{spaces-morphisms-lemma-base-change-surjective}
et
Descent on Spaces,
Lemme \ref{spaces-descent-lemma-descending-property-surjective},
\item universellement injectif, voir
Morphisms of Spaces,
Lemme \ref{spaces-morphisms-lemma-base-change-universally-injective}
et
Descent on Spaces,
Lemme \ref{spaces-descent-lemma-descending-property-universally-injective},
\item localement de type fini, voir
Morphisms of Spaces,
Lemme \ref{spaces-morphisms-lemma-base-change-finite-type}
et
Descent on Spaces,
Lemme \ref{spaces-descent-lemma-descending-property-locally-finite-type},
\item localement de présentation finie, voir
Morphisms of Spaces,
Lemme \ref{spaces-morphisms-lemma-base-change-finite-presentation}
et
Descent on Spaces, Lemme
\ref{spaces-descent-lemma-descending-property-locally-finite-presentation},
\item de type fini, voir
Morphisms of Spaces,
Lemme \ref{spaces-morphisms-lemma-base-change-finite-type}
et
Descent on Spaces,
Lemme \ref{spaces-descent-lemma-descending-property-finite-type},
\item de présentation finie, voir
Morphisms of Spaces,
Lemme \ref{spaces-morphisms-lemma-base-change-finite-presentation}
et
Descent on Spaces, Lemme
\ref{spaces-descent-lemma-descending-property-finite-presentation},
\item plat, voir
Morphisms of Spaces,
Lemme \ref{spaces-morphisms-lemma-base-change-flat}
et
Descent on Spaces,
Lemme \ref{spaces-descent-lemma-descending-property-flat},
\item une immersion ouverte, voir
Morphisms of Spaces,
Section \ref{spaces-morphisms-section-immersions}
et
Descent on Spaces,
Lemme \ref{spaces-descent-lemma-descending-property-open-immersion},
\item un isomorphisme, voir
Descent on Spaces,
Lemme \ref{spaces-descent-lemma-descending-property-isomorphism},
\item affine, voir
Morphisms of Spaces,
Lemme \ref{spaces-morphisms-lemma-base-change-affine}
et
Descent on Spaces,
Lemme \ref{spaces-descent-lemma-descending-property-affine},
\item une immersion fermée, voir
Morphisms of Spaces, Section \ref{spaces-morphisms-section-immersions}
et
Descent on Spaces,
Lemme \ref{spaces-descent-lemma-descending-property-closed-immersion},
\item séparé, voir
Morphisms of Spaces,
Lemme \ref{spaces-morphisms-lemma-base-change-separated}
et
Descent on Spaces,
Lemme \ref{spaces-descent-lemma-descending-property-separated},
\item propre, voir
Morphisms of Spaces,
Lemme \ref{spaces-morphisms-lemma-base-change-proper}
et
Descent on Spaces,
Lemme \ref{spaces-descent-lemma-descending-property-proper},
\item quasi-affine, voir
Morphisms of Spaces,
Lemme \ref{spaces-morphisms-lemma-base-change-quasi-affine}
et
Descent on Spaces,
Lemme \ref{spaces-descent-lemma-descending-property-quasi-affine},
\item entier, voir
Morphisms of Spaces,
Lemme \ref{spaces-morphisms-lemma-base-change-integral}
et
Descent on Spaces,
Lemme \ref{spaces-descent-lemma-descending-property-integral},
\item fini, voir
Morphisms of Spaces,
Lemme \ref{spaces-morphisms-lemma-base-change-integral}
et
Descent on Spaces,
Lemme \ref{spaces-descent-lemma-descending-property-finite},
\item quasi-fini (localement), voir
Morphisms of Spaces,
Lemme \ref{spaces-morphisms-lemma-base-change-quasi-finite}
et
Descent on Spaces,
Lemme \ref{spaces-descent-lemma-descending-property-quasi-finite},
\item syntomique, voir
Morphisms of Spaces,
Lemme \ref{spaces-morphisms-lemma-base-change-syntomic}
et
Descent on Spaces,
Lemme \ref{spaces-descent-lemma-descending-property-syntomic},
\item lisse, voir
Morphisms of Spaces,
Lemme \ref{spaces-morphisms-lemma-base-change-smooth}
et
Descent on Spaces,
Lemme \ref{spaces-descent-lemma-descending-property-smooth},
\item non ramifié, voir
Morphisms of Spaces,
Lemme \ref{spaces-morphisms-lemma-base-change-unramified}
et
Descent on Spaces,
Lemme \ref{spaces-descent-lemma-descending-property-unramified},
\item \'etale, voir
Morphisms of Spaces,
Lemme \ref{spaces-morphisms-lemma-base-change-etale}
et
Descent on Spaces,
Lemme \ref{spaces-descent-lemma-descending-property-etale},
\item fini localement libre, voir
Morphisms of Spaces,
Lemme \ref{spaces-morphisms-lemma-base-change-finite-locally-free}
et
Descent on Spaces,
Lemme \ref{spaces-descent-lemma-descending-property-finite-locally-free},
\item un monomorphisme, voir
Morphisms of Spaces,
Lemme \ref{spaces-morphisms-lemma-base-change-monomorphism}
et
Descent on Spaces,
Lemme \ref{spaces-descent-lemma-descending-property-monomorphism},
\item une immersion, voir
Morphisms of Spaces, Section \ref{spaces-morphisms-section-immersions}
et
Descent on Spaces,
Lemme \ref{spaces-descent-lemma-descending-fppf-property-immersion},
\item localement séparé, voir
Morphisms of Spaces,
Lemme \ref{spaces-morphisms-lemma-base-change-separated}
et
Descent on Spaces,
Lemme \ref{spaces-descent-lemma-descending-fppf-property-locally-separated},
\end{enumerate}

\begin{lemma}
\label{lemma-property-spaces-too}
Soit $P$ une propriété des morphismes d'espaces algébriques comme ci-dessus.
Soit $f : \mathcal{X} \to \mathcal{Y}$ un morphisme de champs algébriques
représentable par des espaces algébriques. Les assertions suivantes sont équivalentes :
\begin{enumerate}
\item $f$ possède $P$,
\item pour tout espace algébrique $Z$ et tout morphisme $Z \to \mathcal{Y}$,
le morphisme $Z \times_\mathcal{Y} \mathcal{X} \to Z$ possède $P$.
\end{enumerate}
\end{lemma}

\begin{proof}
L'implication (2) $\Rightarrow$ (1) est immédiate. Supposons (1).
Soit $Z \to \mathcal{Y}$ comme en (2). Choisissons un schéma $U$ et un
morphisme \'etale surjectif $U \to Z$. Par hypothèse, le morphisme
$U \times_\mathcal{Y} \mathcal{X} \to U$ possède $P$. Mais le diagramme
$$
\xymatrix{
U \times_\mathcal{Y} \mathcal{X} \ar[d] \ar[r] &
Z \times_\mathcal{Y} \mathcal{X} \ar[d] \\
U \ar[r] & Z
}
$$
est cartésien; la flèche verticale de droite possède donc $P$, puisque
$\{U \to Z\}$ est un recouvrement fppf.
\end{proof}

\noindent
Le lemme suivant nous dit qu'il suffit de vérifier $P$
après un changement de base par un morphisme surjectif, plat et localement
de présentation finie.

\begin{lemma}
\label{lemma-check-property-covering}
Soit $P$ une propriété des morphismes d'espaces algébriques comme ci-dessus.
Soit $f : \mathcal{X} \to \mathcal{Y}$ un morphisme de champs algébriques
représentable par des espaces algébriques.
Soit $W$ un espace algébrique et soit $W \to \mathcal{Y}$ un morphisme surjectif,
localement de présentation finie et plat.
Posons $V = W \times_\mathcal{Y} \mathcal{X}$. Alors
$$
(f\text{ possède }P) \Leftrightarrow (\text{la projection }V \to W\text{ possède }P).
$$
\end{lemma}

\begin{proof}
L'implication de gauche à droite résulte du
Lemme \ref{lemma-property-spaces-too}.
Supposons que $V \to W$ possède $P$. Soit $T$ un schéma et soit
$T \to \mathcal{Y}$ un morphisme. Considérons le diagramme commutatif
$$
\xymatrix{
T \times_\mathcal{Y} \mathcal{X} \ar[d] &
T \times_\mathcal{Y} V \ar[d] \ar[l] \ar[r] &
V \ar[d] \\
T & T \times_\mathcal{Y} W \ar[l] \ar[r] & W
}
$$
d'espaces algébriques. Les carrés sont cartésiens.
Le morphisme inférieur de gauche est un morphisme surjectif et plat, localement de
présentation finie; $\{T \times_\mathcal{Y} V \to T\}$ est donc un
recouvrement fppf. Par conséquent, le fait que la flèche verticale de droite possède la propriété
$P$ implique que la flèche verticale de gauche possède la propriété $P$.
\end{proof}

\begin{lemma}
\label{lemma-check-property-weak-covering}
Soit $P$ une propriété des morphismes d'espaces algébriques comme ci-dessus.
Soit $f : \mathcal{X} \to \mathcal{Y}$ un morphisme de champs algébriques
représentable par des espaces algébriques.
Soit $\mathcal{Z} \to \mathcal{Y}$ un morphisme de champs algébriques qui
est représentable par des espaces algébriques, surjectif, plat et
localement de présentation finie.
Posons $\mathcal{W} = \mathcal{Z} \times_\mathcal{Y} \mathcal{X}$. Alors
$$
(f\text{ possède }P) \Leftrightarrow
(\text{la projection }\mathcal{W} \to \mathcal{Z}\text{ possède }P).
$$
\end{lemma}

\begin{proof}
Choisissons un espace algébrique $W$ et un morphisme
$W \to \mathcal{Z}$ qui soit surjectif, plat et localement de présentation
finie. D'après la discussion ci-dessus, le composé
$W \to \mathcal{Y}$ est aussi surjectif, plat et localement de présentation
finie. Notons
$V = W \times_\mathcal{Z} \mathcal{W} = V \times_\mathcal{Y} \mathcal{X}$.
Par le
Lemme \ref{lemma-check-property-covering},
nous voyons que $f$ possède $\mathcal{P}$ si et seulement si $V \to W$ la possède,
et que $\mathcal{W} \to \mathcal{Z}$ possède $\mathcal{P}$ si et seulement
si $V \to W$ la possède. Le lemme en résulte.
\end{proof}

\begin{lemma}
\label{lemma-check-property-after-precomposing}
Soit $P$ une propriété des morphismes d'espaces algébriques comme ci-dessus.
Soit $\tau \in \{\etale, smooth, syntomic, fppf\}$.
Soient $\mathcal{X} \to \mathcal{Y}$ et $\mathcal{Y} \to \mathcal{Z}$
des morphismes de champs algébriques représentables par des espaces algébriques.
Supposons que
\begin{enumerate}
\item $\mathcal{X} \to \mathcal{Y}$ soit surjectif et
\'etale, lisse, syntomique, ou plat et localement de présentation finie,
\item le composé possède $P$, et
\item $P$ soit locale sur la source pour la topologie $\tau$.
\end{enumerate}
Alors $\mathcal{Y} \to \mathcal{Z}$ possède la propriété $P$.
\end{lemma}

\begin{proof}
Soit $Z$ un schéma et soit $Z \to \mathcal{Z}$ un morphisme.
Posons $X = \mathcal{X} \times_\mathcal{Z} Z$,
$Y = \mathcal{Y} \times_\mathcal{Z} Z$. D'après (1), $\{X \to Y\}$
est un recouvrement pour la topologie $\tau$ d'espaces algébriques, et, d'après (2), $X \to Z$ possède la propriété
$P$. D'après (3), il en résulte que $Y \to Z$ possède la propriété $P$,
ce qui conclut la démonstration.
\end{proof}

\begin{lemma}
\label{lemma-representable-in-terms-presentations}
Soit $g : \mathcal{X}' \to \mathcal{X}$ un morphisme de champs algébriques
qui soit représentable par des espaces algébriques. Soit $[U/R] \to \mathcal{X}$
une présentation. Posons $U' = U \times_\mathcal{X} \mathcal{X}'$,
et $R' = R \times_\mathcal{X} \mathcal{X}'$.
Il existe alors un groupoïde en espaces algébriques de la forme
$(U', R', s', t', c')$, une présentation $[U'/R'] \to \mathcal{X}'$,
et le diagramme
$$
\xymatrix{
[U'/R'] \ar[d]_{[\text{pr}]} \ar[r] & \mathcal{X}' \ar[d]^g \\
[U/R] \ar[r] & \mathcal{X}
}
$$
est $2$-commutatif, où le morphisme $[\text{pr}]$ provient d'un
morphisme de groupoïdes
$\text{pr} : (U', R', s', t', c') \to (U, R, s, t, c)$.
\end{lemma}

\begin{proof}
Puisque $U \to \mathcal{Y}$ est surjectif et lisse, voir
Algebraic Stacks,
Lemme \ref{algebraic-lemma-smooth-quotient-smooth-presentation},
le changement de base $U' \to \mathcal{X}'$ est aussi surjectif et lisse.
Par suite, d'après
Algebraic Stacks,
Lemme \ref{algebraic-lemma-stack-presentation},
il suffit de montrer que $R' = U' \times_{\mathcal{X}'} U'$ afin
d'obtenir un groupoïde lisse $(U', R', s', t', c')$ et une présentation
$[U'/R'] \to \mathcal{X}'$.
Comme $R = V \times_\mathcal{Y} V$ (voir
Groupoids in Spaces,
Lemme \ref{spaces-groupoids-lemma-quotient-stack-2-cartesian}),
cela résulte de
$$
R' =
U \times_\mathcal{X} U \times_\mathcal{X} \mathcal{X}' =
(U \times_\mathcal{X} \mathcal{X}')
\times_{\mathcal{X}'}
(U \times_\mathcal{X} \mathcal{X}')
$$
voir
Categories, Lemmes \ref{categories-lemma-associativity-2-fibre-product} et
\ref{categories-lemma-2-fibre-product-erase-factor}.
Il est clair que les morphismes de projection $U' \to U$ et $R' \to R$ donnent le
morphisme de groupoïdes désiré
$\text{pr} : (U', R', s', t', c') \to (U, R, s, t, c)$.
On obtient donc le morphisme $[\text{pr}]$ des champs quotients grâce à
Groupoids in Spaces,
Lemme \ref{spaces-groupoids-lemma-quotient-stack-functorial}.

\medskip\noindent
Il nous reste à montrer que le diagramme est $2$-commutatif.
Il est clair que le diagramme
$$
\xymatrix{
U' \ar[d]_{\text{pr}_U} \ar[r]_{f'} & \mathcal{X}' \ar[d]^g \\
U \ar[r]^f & \mathcal{X}
}
$$
est $2$-commutatif, où $\text{pr}_U : U' \to U$ est la projection.
Il existe une $2$-flèche canonique
$\tau : f \circ t \to f \circ s$ dans $\Mor(R, \mathcal{X})$
provenant de $R = U \times_\mathcal{X} U$, $t = \text{pr}_0$, et
$s = \text{pr}_1$. En utilisant l'isomorphisme
$R' \to U' \times_{\mathcal{X}'} U'$, nous obtenons de même un isomorphisme
$\tau' : f' \circ t' \to f' \circ s'$. Remarquons que
$g \circ f' \circ t' = f \circ t \circ \text{pr}_R$ et
$g \circ f' \circ s' = f \circ s \circ \text{pr}_R$, où
$\text{pr}_R : R' \to R$ est la projection. Il est donc légitime de se demander
si
\begin{equation}
\label{equation-verify}
\tau \star \text{id}_{\text{pr}_R} = \text{id}_g \star \tau'.
\end{equation}
Nous avançons maintenant deux assertions : (1) si l'équation (\ref{equation-verify}) est satisfaite,
alors le diagramme est $2$-commutatif, et (2) l'équation (\ref{equation-verify}) est satisfaite.
Nous omettons la démonstration de ces deux assertions. Indications : la partie (1) résulte de la
construction de $f = f_{can}$ et $f' = f'_{can}$ dans
Algebraic Stacks, Lemme \ref{algebraic-lemma-map-space-into-stack}.
La partie (2) résulte d'un examen attentif des définitions.
\end{proof}

\begin{remark}
\label{remark-representable-over}
Soit $\mathcal{Y}$ un champ algébrique. Considérons la
$2$-catégorie suivante :
\begin{enumerate}
\item un objet est un morphisme $f : \mathcal{X} \to \mathcal{Y}$
représentable par des espaces algébriques,
\item un $1$-morphisme
$(g, \beta) :
(f_1 : \mathcal{X}_1 \to \mathcal{Y})
\to
(f_2 : \mathcal{X}_2 \to \mathcal{Y})$
consiste en un morphisme $g : \mathcal{X}_1 \to \mathcal{X}_2$ et un
$2$-morphisme $\beta : f_1 \to f_2 \circ g$, et
\item un $2$-morphisme entre
$(g, \beta), (g', \beta') :
(f_1 : \mathcal{X}_1 \to \mathcal{Y})
\to
(f_2 : \mathcal{X}_2 \to \mathcal{Y})$
est un $2$-morphisme $\alpha : g \to g'$ tel que
$(\text{id}_{f_2} \star \alpha) \circ \beta = \beta'$.
\end{enumerate}
Notons cette $2$-catégorie $\textit{Spaces}/\mathcal{Y}$, par
analogie avec la notation de
Topologies on Spaces, Section \ref{spaces-topologies-section-procedure}.
Nous affirmons que, dans cette $2$-catégorie, les catégories de morphismes
$$
\Mor_{\textit{Spaces}/\mathcal{Y}}(
(f_1 : \mathcal{X}_1 \to \mathcal{Y}),
(f_2 : \mathcal{X}_2 \to \mathcal{Y}))
$$
sont toutes des sétoïdes. En effet, un $2$-morphisme $\alpha$ est une règle qui, à tout
objet $x_1$ de $\mathcal{X}_1$, associe un isomorphisme
$\alpha_{x_1} : g(x_1) \longrightarrow g'(x_1)$
dans la catégorie fibre correspondante de $\mathcal{X}_2$, de telle sorte que le diagramme
$$
\xymatrix{
& f_2(x_1) \ar[ld]_{\beta_{x_1}} \ar[rd]^{\beta'_{x_1}} \\
f_2(g(x_1)) \ar[rr]^{f_2(\alpha_{x_1})} & &
f_2(g'(x_1))
}
$$
commute. Mais, puisque $f_2$ est fidèle (voir
Algebraic Stacks,
Lemme \ref{algebraic-lemma-characterize-representable-by-algebraic-spaces}),
cela signifie que, si $\alpha_{x_1}$ existe, il est unique ! Autrement dit, la
$2$-catégorie $\textit{Spaces}/\mathcal{Y}$ est très proche d'être
une catégorie. Plus précisément, si nous remplaçons les $1$-morphismes par les classes
d'isomorphisme de $1$-morphismes, nous obtenons une catégorie. Nous effectuerons souvent
ce remplacement sans autre mention.
\end{remark}




\section{Points des champs algébriques}
\label{section-points}

\noindent
Soit $\mathcal{X}$ un champ algébrique. Soient $K, L$ deux corps,
et soient $p : \Spec(K) \to \mathcal{X}$ et
$q : \Spec(L) \to \mathcal{X}$ des morphismes.
Nous disons que $p$ et $q$ sont {\it équivalents} s'il existe un
corps $\Omega$ et un diagramme $2$-commutatif
$$
\xymatrix{
\Spec(\Omega) \ar[r] \ar[d] &
\Spec(L) \ar[d]^q \\
\Spec(K) \ar[r]^p &
\mathcal{X}.
}
$$

\begin{lemma}
\label{lemma-equivalence}
La notion ci-dessus définit bien une relation d'équivalence sur les
morphismes des spectres de corps vers le champ algébrique $\mathcal{X}$.
\end{lemma}

\begin{proof}
Il est clair que la relation est réflexive et symétrique.
Il reste donc à montrer qu'elle est transitive. Cela revient
à ceci : étant donné un diagramme
$$
\xymatrix{
\Spec(\Omega) \ar[r]_b \ar[d]_a &
\Spec(L) \ar[d]^q & \Spec(\Omega') \ar[l]^{b'} \ar[d]^{a'} \\
\Spec(K) \ar[r]^p &
\mathcal{X} &
\Spec(K') \ar[l]_{p'}
}
$$
dont les deux carrés sont $2$-commutatifs, nous devons montrer que $p$ est équivalent à
$p'$. Par le lemme de $2$-Yoneda (voir
Algebraic Stacks, Section \ref{algebraic-section-2-yoneda}),
les morphismes $p$, $p'$ et $q$ sont donnés par des objets
$x$, $x'$ et $y$ dans les catégories fibres de $\mathcal{X}$ au-dessus de
$\Spec(K)$, $\Spec(K')$ et $\Spec(L)$. La
$2$-commutativité des carrés signifie qu'il existe des isomorphismes
$\alpha : a^*x \to b^*y$ et $\alpha' : (a')^*x' \to (b')^*y$
dans les catégories fibres
de $\mathcal{X}$ au-dessus de $\Spec(\Omega)$ et $\Spec(\Omega')$.
Choisissons un corps $\Omega''$ et des plongements
$\Omega \to \Omega''$ et $\Omega' \to \Omega''$ coïncidant sur $L$.
Nous pouvons alors prolonger le diagramme ci-dessus en
$$
\xymatrix{
& \Spec(\Omega'') \ar[ld]_c \ar[d]^{q'} \ar[rd]^{c'} \\
\Spec(\Omega) \ar[r]_b \ar[d]_a &
\Spec(L) \ar[d]^q & \Spec(\Omega') \ar[l]^{b'} \ar[d]^{a'} \\
\Spec(K) \ar[r]^p &
\mathcal{X} &
\Spec(K') \ar[l]_{p'}
}
$$
avec des triangles commutatifs, et
$$
(q')^*(\alpha')^{-1} \circ (q')^*\alpha :
(a \circ c)^*x
\longrightarrow
(a' \circ c')^*x'
$$
est un isomorphisme dans la catégorie fibre au-dessus de $\Spec(\Omega'')$.
Ainsi, $p$ est équivalent à $p'$, comme souhaité.
\end{proof}

\begin{definition}
\label{definition-points}
Soit $\mathcal{X}$ un champ algébrique.
Un {\it point} de $\mathcal{X}$ est une classe d'équivalence de morphismes
de spectres de corps vers $\mathcal{X}$.
L'ensemble des points de $\mathcal{X}$ est noté $|\mathcal{X}|$.
\end{definition}

\noindent
Cela concorde avec notre définition des points des espaces algébriques, voir
Properties of Spaces, Définition \ref{spaces-properties-definition-points}.
De plus, pour un schéma, nous retrouvons la notion usuelle de point, voir
Properties of Spaces, Lemme \ref{spaces-properties-lemma-scheme-points}.
Si $f : \mathcal{X} \to \mathcal{Y}$ est un morphisme de champs algébriques,
il existe une application induite $|f| : |\mathcal{X}| \to |\mathcal{Y}|$ qui
envoie un représentant $x : \Spec(K) \to \mathcal{X}$ sur le
représentant $f \circ x : \Spec(K) \to \mathcal{Y}$. Celle-ci est
bien définie : en effet, lorsqu'ils sont $2$-isomorphes, des $1$-morphismes restent $2$-isomorphes après
pré- ou postcomposition par un $1$-morphisme, car on peut les pré- ou postcomposer horizontalement
par l'identité du $1$-morphisme donné. Cela vaut
dans toute $(2, 1)$-catégorie (stricte). Si
$$
\xymatrix{
\mathcal{X} \ar[d] \ar[r] & \mathcal{Y} \ar[d] \\
\mathcal{W} \ar[r] & \mathcal{Z}
}
$$
est un diagramme $2$-commutatif de champs algébriques, alors le diagramme
d'ensembles
$$
\xymatrix{
|\mathcal{X}| \ar[d] \ar[r] & |\mathcal{Y}| \ar[d] \\
|\mathcal{W}| \ar[r] & |\mathcal{Z}|
}
$$
est commutatif.
En particulier, si $\mathcal{X} \to \mathcal{Y}$ est une équivalence,
alors $|\mathcal{X}| \to |\mathcal{Y}|$ est une bijection.

\begin{lemma}
\label{lemma-points-cartesian}
Soit
$$
\xymatrix{
\mathcal{Z} \times_\mathcal{Y} \mathcal{X} \ar[r] \ar[d] &
\mathcal{X} \ar[d] \\
\mathcal{Z} \ar[r] & \mathcal{Y}
}
$$
un produit fibré de champs algébriques. Alors l'application d'ensembles
de points
$$
|\mathcal{Z} \times_\mathcal{Y} \mathcal{X}|
\longrightarrow
|\mathcal{Z}| \times_{|\mathcal{Y}|} |\mathcal{X}|
$$
est surjective.
\end{lemma}

\begin{proof}
En effet, supposons donnés des corps $K$, $L$ et des morphismes
$\Spec(K) \to \mathcal{X}$, $\Spec(L) \to \mathcal{Z}$.
L'hypothèse qu'ils coïncident comme éléments de $|\mathcal{Y}|$ signifie alors
qu'il existe une extension commune $M/K$ et $M/L$
telle que
$\Spec(M) \to \Spec(K) \to \mathcal{X} \to \mathcal{Y}$ et
$\Spec(M) \to \Spec(L) \to \mathcal{Z} \to \mathcal{Y}$
soient $2$-isomorphes. Et c'est exactement la condition qui donne un
morphisme $\Spec(M) \to \mathcal{Z} \times_\mathcal{Y} \mathcal{X}$.
\end{proof}

\begin{lemma}
\label{lemma-characterize-surjective}
Soit $f : \mathcal{X} \to \mathcal{Y}$ un morphisme de champs algébriques
représentable par des espaces algébriques. Les assertions suivantes sont équivalentes :
\begin{enumerate}
\item $|f| : |\mathcal{X}| \to |\mathcal{Y}|$ est surjective, et
\item $f$ est surjectif
(au sens de la Section \ref{section-properties-morphisms}).
\end{enumerate}
\end{lemma}

\begin{proof}
Supposons (1). Soit $T \to \mathcal{Y}$ un morphisme dont la source est un schéma.
Pour démontrer (2), nous devons montrer que le morphisme d'espaces algébriques
$T \times_\mathcal{Y} \mathcal{X} \to T$ est surjectif. D'après
Morphisms of Spaces, Définition \ref{spaces-morphisms-definition-surjective},
cela signifie que nous devons montrer que
$|T \times_\mathcal{Y} \mathcal{X}| \to |T|$ est surjective.
En appliquant le
Lemme \ref{lemma-points-cartesian},
nous voyons que cela résulte de (1).

\medskip\noindent
Réciproquement, supposons (2). Soit $y : \Spec(K) \to \mathcal{Y}$ un
morphisme du spectre d'un corps vers $\mathcal{Y}$. Par hypothèse, le
morphisme
$\Spec(K) \times_{y, \mathcal{Y}} \mathcal{X} \to \Spec(K)$
d'espaces algébriques est surjectif. D'après
Morphisms of Spaces, Définition \ref{spaces-morphisms-definition-surjective},
cela signifie qu'il existe une extension de corps
$K'/K$ et un morphisme
$\Spec(K') \to \Spec(K) \times_{y, \mathcal{Y}} \mathcal{X}$
tels que le carré de gauche du diagramme
$$
\xymatrix{
\Spec(K') \ar[r] \ar[d] &
\Spec(K) \times_{y, \mathcal{Y}} \mathcal{X} \ar[d] \ar[r] &
\mathcal{X} \ar[d]
\\
\Spec(K) \ar@{=}[r] &
\Spec(K) \ar[r]^-y &
\mathcal{Y}
}
$$
soit commutatif. Cela montre que $|X| \to |\mathcal{Y}|$ est surjective.
\end{proof}

\noindent
Voici un lemme qui explique comment calculer l'ensemble des points à partir
d'une présentation.

\begin{lemma}
\label{lemma-points-presentation}
Soit $\mathcal{X}$ un champ algébrique.
Soit $\mathcal{X} = [U/R]$ une présentation de $\mathcal{X}$, voir
Algebraic Stacks, Définition \ref{algebraic-definition-presentation}.
Alors l'image de $|R| \to |U| \times |U|$ est une relation d'équivalence
et $|\mathcal{X}|$ est le quotient de $|U|$ par cette relation d'équivalence.
\end{lemma}

\begin{proof}
L'hypothèse signifie que nous avons un groupoïde lisse $(U, R, s, t, c)$
en espaces algébriques, et une équivalence $f : [U/R] \to \mathcal{X}$.
Nous pouvons supposer que $\mathcal{X} = [U/R]$.
Le morphisme induit $p : U \to \mathcal{X}$ est lisse et surjectif, voir
Algebraic Stacks,
Lemme \ref{algebraic-lemma-smooth-quotient-smooth-presentation}.
Par conséquent, $|U| \to |\mathcal{X}|$ est surjective d'après le
Lemme \ref{lemma-characterize-surjective}.
Remarquons que $R = U \times_\mathcal{X} U$, voir
Groupoids in Spaces,
Lemme \ref{spaces-groupoids-lemma-quotient-stack-2-cartesian}.
Le
Lemme \ref{lemma-points-cartesian}
implique donc que l'application
$$
|R| \longrightarrow |U| \times_{|\mathcal{X}|} |U|
$$
est surjective. L'image de $|R| \to |U| \times |U|$ est donc
exactement l'ensemble des couples $(u_1, u_2) \in |U| \times |U|$
tels que $u_1$ et $u_2$ aient la même image dans $|\mathcal{X}|$.
En combinant ces deux assertions, nous obtenons le résultat du lemme.
\end{proof}

\begin{remark}
\label{remark-more-general-presentation}
Le résultat du
Lemme \ref{lemma-points-presentation}
se généralise comme suit.
Soit $\mathcal{X}$ un champ algébrique.
Soit $U$ un espace algébrique et soit $f : U \to \mathcal{X}$ un morphisme surjectif
(ce qui a un sens d'après la
Section \ref{section-properties-morphisms}).
Soit $R = U \times_\mathcal{X} U$, soit $(U, R, s, t, c)$ le groupoïde
en espaces algébriques et soit $f_{can} : [U/R] \to \mathcal{X}$ le
morphisme canonique construit dans
Algebraic Stacks, Lemme \ref{algebraic-lemma-map-space-into-stack}.
Alors l'image de $|R| \to |U| \times |U|$ est une relation d'équivalence
et $|\mathcal{X}| = |U|/|R|$. La démonstration du
Lemme \ref{lemma-points-presentation}
s'applique sans changement. (Bien entendu, en général, $[U/R]$ n'est pas un champ
algébrique et, en général, $f_{can}$ n'est pas un isomorphisme.)
\end{remark}

\begin{lemma}
\label{lemma-topology-points}
Il existe une unique topologie sur les ensembles de points
des champs algébriques qui possède les propriétés suivantes :
\begin{enumerate}
\item pour tout morphisme de champs algébriques $\mathcal{X} \to \mathcal{Y}$,
l'application $|\mathcal{X}| \to |\mathcal{Y}|$ est continue, et
\item pour tout morphisme $U \to \mathcal{X}$ qui est plat et localement
de présentation finie, où $U$ est un espace algébrique,
l'application d'espaces topologiques $|U| \to |\mathcal{X}|$ est continue et ouverte.
\end{enumerate}
\end{lemma}

\begin{proof}
Choisissons un morphisme $p : U \to \mathcal{X}$ qui soit
surjectif, plat et localement de présentation finie,
où $U$ est un espace algébrique. Un tel morphisme existe par définition d'un champ
algébrique, puisqu'un morphisme lisse est plat et localement de présentation finie
(voir
Morphisms of Spaces,
Lemmes \ref{spaces-morphisms-lemma-smooth-locally-finite-presentation} et
\ref{spaces-morphisms-lemma-smooth-flat}).
Nous définissons une topologie sur $|\mathcal{X}|$ par la règle suivante :
$W \subset |\mathcal{X}|$ est ouvert si et seulement si $|p|^{-1}(W)$ est ouvert
dans $|U|$. Pour montrer que cette définition est indépendante du choix de $p$, soit
$p' : U' \to \mathcal{X}$ un autre morphisme surjectif, plat et
localement de présentation finie d'un espace algébrique vers
$\mathcal{X}$. Posons $U'' = U \times_\mathcal{X} U'$,
de sorte que nous avons un diagramme $2$-commutatif
$$
\xymatrix{
U'' \ar[r] \ar[d] & U' \ar[d] \\
U \ar[r] & \mathcal{X}
}
$$
Puisque $U \to \mathcal{X}$ et $U' \to \mathcal{X}$ sont surjectifs, plats et
localement de présentation finie, nous voyons que $U'' \to U'$ et $U'' \to U$
sont surjectifs, plats et localement de présentation finie, voir
Lemme \ref{lemma-property-spaces-too}.
Par conséquent, les applications $|U''| \to |U'|$ et $|U''| \to |U|$ sont continues, ouvertes
et surjectives, voir
Morphisms of Spaces,
Définition \ref{spaces-morphisms-definition-surjective} et
Lemme \ref{spaces-morphisms-lemma-fppf-open}.
Cela implique clairement que notre définition est indépendante du choix
de $p : U \to \mathcal{X}$.

\medskip\noindent
Soit $f : \mathcal{X} \to \mathcal{Y}$ un morphisme de champs algébriques.
D'après
Algebraic Stacks, Lemme \ref{algebraic-lemma-lift-morphism-presentations},
nous pouvons trouver un diagramme $2$-commutatif
$$
\xymatrix{
U \ar[d]_x \ar[r]_a & V \ar[d]^y \\
\mathcal{X} \ar[r]^f & \mathcal{Y}
}
$$
dont les flèches verticales sont lisses et surjectives.
Considérons le diagramme commutatif associé
$$
\xymatrix{
|U| \ar[d]_{|x|} \ar[r]_{|a|} & |V| \ar[d]^{|y|} \\
|\mathcal{X}| \ar[r]^{|f|} & |\mathcal{Y}|
}
$$
d'ensembles. Si $W \subset |\mathcal{Y}|$ est ouvert, alors, par la définition
ci-dessus, cela signifie exactement que $|y|^{-1}(W)$ est ouvert dans $|V|$. Puisque
$|a|$ est continue, nous en concluons que
$|a|^{-1}|y|^{-1}(W) = |x|^{-1}|f|^{-1}(W)$ est ouvert dans $|W|$, ce qui signifie,
par définition, que $|f|^{-1}(W)$ est ouvert dans $|\mathcal{X}|$.
Ainsi, $|f|$ est continue.

\medskip\noindent
Enfin, nous devons montrer que si $U$ est un espace algébrique et si
$U \to \mathcal{X}$ est plat et localement de présentation finie, alors
$|U| \to |\mathcal{X}|$ est ouverte. Soit $V \to \mathcal{X}$ un morphisme surjectif,
plat et localement de présentation finie, où $V$ est un espace algébrique.
Considérons le diagramme commutatif
$$
\xymatrix{
|U \times_\mathcal{X} V| \ar[r]_e \ar[rd]_f &
|U| \times_{|\mathcal{X}|} |V| \ar[d]_c \ar[r]_d &
|V| \ar[d]^b \\
& |U| \ar[r]^a & |\mathcal{X}|
}
$$
Le morphisme $U \times_\mathcal{X} V \to U$ est surjectif, c'est-à-dire que
$f : |U \times_\mathcal{X} V| \to |U|$ est surjective.
La flèche horizontale supérieure de gauche est surjective, voir
Lemme \ref{lemma-points-cartesian}.
Le morphisme $U \times_\mathcal{X} V \to V$ est plat et localement de présentation
finie; $d \circ e : |U \times_\mathcal{X} V| \to |V|$ est donc ouverte,
voir
Morphisms of Spaces, Lemme \ref{spaces-morphisms-lemma-fppf-open}.
Choisissons un ouvert $W \subset |U|$. Les propriétés ci-dessus impliquent que
$b^{-1}(a(W)) = (d \circ e)(f^{-1}(W))$ est ouvert, ce qui, par construction, signifie
que $a(W)$ est ouvert, comme souhaité.
\end{proof}

\begin{definition}
\label{definition-topological-space}
Soit $\mathcal{X}$ un champ algébrique.
L'{\it espace topologique} sous-jacent à $\mathcal{X}$ est l'ensemble des points
$|\mathcal{X}|$ muni de la topologie construite dans le
Lemme \ref{lemma-topology-points}.
\end{definition}

\noindent
Cette définition n'entre pas en conflit avec la topologie déjà existante
sur $|\mathcal{X}|$ lorsque $\mathcal{X}$ est un espace algébrique.

\begin{lemma}
\label{lemma-space-locally-quasi-compact}
Soit $\mathcal{X}$ un champ algébrique.
Tout point de $|\mathcal{X}|$ possède un système fondamental de
voisinages ouverts quasi-compacts.
En particulier, $|\mathcal{X}|$ est localement quasi-compact au sens de
Topology, Définition \ref{topology-definition-locally-quasi-compact}.
\end{lemma}

\begin{proof}
Cela résulte formellement du fait qu'il existe un schéma
$U$ et une application surjective, ouverte et continue $U \to |\mathcal{X}|$
d'espaces topologiques. En effet, si $U \to \mathcal{X}$ est surjectif et
lisse, alors le
Lemme \ref{lemma-topology-points}
garantit que $|U| \to |\mathcal{X}|$ est continue, surjective et ouverte.
\end{proof}











\section{Morphismes surjectifs}
\label{section-surjective}

\noindent
Soit $f : \mathcal{X} \to \mathcal{Y}$ un morphisme de champs algébriques
représentable par des espaces algébriques. Dans la
Section \ref{section-properties-morphisms},
nous avons déjà défini ce que signifie pour $f$ d'être surjectif. Dans le
Lemme \ref{lemma-characterize-surjective},
nous avons vu que cela équivaut à demander que
$|f| : |\mathcal{X}| \to |\mathcal{Y}|$ soit surjective.
Cela ouvre la voie à la définition suivante.

\begin{definition}
\label{definition-surjective}
Soit $f : \mathcal{X} \to \mathcal{Y}$ un morphisme de champs algébriques.
Nous disons que $f$ est {\it surjectif} si l'application
$|f| : |\mathcal{X}| \to |\mathcal{Y}|$ des espaces topologiques associés
est surjective.
\end{definition}

\noindent
Voici quelques lemmes.

\begin{lemma}
\label{lemma-composition-surjective}
Le composé de morphismes surjectifs est surjectif.
\end{lemma}

\begin{proof}
Omis.
\end{proof}

\begin{lemma}
\label{lemma-base-change-surjective}
Tout changement de base d'un morphisme surjectif est surjectif.
\end{lemma}

\begin{proof}
Omis. Indication : utiliser le
Lemme \ref{lemma-points-cartesian}.
\end{proof}

\begin{lemma}
\label{lemma-descent-surjective}
Soit $f : \mathcal{X} \to \mathcal{Y}$ un morphisme de champs algébriques.
Soit $\mathcal{Y}' \to \mathcal{Y}$ un morphisme surjectif de champs
algébriques. Si le changement de base $f' : \mathcal{Y}' \times_\mathcal{Y} \mathcal{X}
\to \mathcal{Y}'$ de $f$ est surjectif, alors $f$ est surjectif.
\end{lemma}

\begin{proof}
Cela résulte immédiatement du
Lemme \ref{lemma-points-cartesian}.
\end{proof}

\begin{lemma}
\label{lemma-surjective-permanence}
Soient $\mathcal{X} \to \mathcal{Y} \to \mathcal{Z}$ des morphismes de
champs algébriques. Si $\mathcal{X} \to \mathcal{Z}$ est surjectif,
alors $\mathcal{Y} \to \mathcal{Z}$ l'est aussi.
\end{lemma}

\begin{proof}
Immédiat.
\end{proof}












\section{Champs algébriques quasi-compacts}
\label{section-quasi-compact}

\noindent
La définition suivante équivaut à la définition pour les espaces
algébriques d'après
Properties of Spaces, Lemme \ref{spaces-properties-lemma-quasi-compact-space}.

\begin{definition}
\label{definition-quasi-compact}
Soit $\mathcal{X}$ un champ algébrique.
Nous disons que $\mathcal{X}$ est {\it quasi-compact}
si et seulement si $|\mathcal{X}|$ est quasi-compact.
\end{definition}

\begin{lemma}
\label{lemma-quasi-compact-stack}
Soit $\mathcal{X}$ un champ algébrique.
Les assertions suivantes sont équivalentes :
\begin{enumerate}
\item $\mathcal{X}$ est quasi-compact,
\item il existe un morphisme lisse surjectif $U \to \mathcal{X}$,
où $U$ est un schéma affine,
\item il existe un morphisme lisse surjectif $U \to \mathcal{X}$,
où $U$ est un schéma quasi-compact,
\item il existe un morphisme lisse surjectif $U \to \mathcal{X}$,
où $U$ est un espace algébrique quasi-compact, et
\item il existe un morphisme surjectif $\mathcal{U} \to \mathcal{X}$
de champs algébriques tel que $\mathcal{U}$ soit quasi-compact.
\end{enumerate}
\end{lemma}

\begin{proof}
Nous utiliserons le
Lemme \ref{lemma-characterize-surjective}.
Supposons que $\mathcal{U}$ et $\mathcal{U} \to \mathcal{X}$ soient comme en (5).
Puisque $|\mathcal{U}| \to |\mathcal{X}|$ est surjective et
continue, nous en concluons que $|\mathcal{X}|$ est quasi-compact.
Ainsi, (5) implique (1). Les implications
(2) $\Rightarrow$ (3) $\Rightarrow$ (4) $\Rightarrow$ (5)
sont immédiates. Supposons (1), c'est-à-dire que $\mathcal{X}$ est quasi-compact, c'est-à-dire que
$|\mathcal{X}|$ est quasi-compact. Choisissons un schéma $U$ et un morphisme lisse
surjectif $U \to \mathcal{X}$. Puisque $|U| \to |\mathcal{X}|$
est ouverte, nous voyons qu'il existe un ouvert quasi-compact $U' \subset U$
tel que $|U'| \to |X|$ soit surjective (et toujours lisse).
Choisissons un recouvrement ouvert affine fini $U' = U_1 \cup \ldots \cup U_n$.
Alors $U_1 \amalg \ldots \amalg U_n \to \mathcal{X}$
est un morphisme lisse surjectif dont la source est un
schéma affine (Schemes, Lemme \ref{schemes-lemma-disjoint-union-affines}).
Ainsi, (2) est satisfaite.
\end{proof}

\begin{lemma}
\label{lemma-finite-disjoint-quasi-compact}
Une réunion disjointe finie de champs algébriques quasi-compacts est
un champ algébrique quasi-compact.
\end{lemma}

\begin{proof}
Cela résulte clairement du fait topologique correspondant.
\end{proof}


\section{Propriétés des champs algébriques définies par des propriétés des schémas}
\label{section-types-properties}

\noindent
Toute propriété des schémas locale pour la topologie lisse donne naissance à une
propriété correspondante des champs algébriques grâce au lemme suivant. Remarquons qu'une
propriété des schémas qui est locale pour la topologie lisse est aussi locale pour la topologie \'etale,
puisque tout recouvrement \'etale est aussi un recouvrement lisse. Ainsi, pour une propriété $P$
des schémas locale pour la topologie lisse, nous savons ce que signifie dire qu'un
espace algébrique possède $P$, voir
Properties of Spaces, Section \ref{spaces-properties-section-types-properties}.

\begin{lemma}
\label{lemma-type-property}
Soit $\mathcal{P}$ une propriété des schémas qui est locale pour la topologie
lisse, voir
Descent, Définition \ref{descent-definition-property-local}.
Soit $\mathcal{X}$ un champ algébrique. Les assertions suivantes sont équivalentes :
\begin{enumerate}
\item il existe un schéma $U$ et un morphisme lisse surjectif
$U \to \mathcal{X}$ tels que le schéma $U$ possède la propriété $\mathcal{P}$,
\item pour tout schéma $U$ et tout morphisme lisse $U \to \mathcal{X}$,
le schéma $U$ possède la propriété $\mathcal{P}$,
\item il existe un espace algébrique $U$ et un morphisme lisse surjectif
$U \to \mathcal{X}$ tels que l'espace algébrique $U$ possède la propriété $\mathcal{P}$, et
\item pour tout espace algébrique $U$ et tout morphisme lisse
$U \to \mathcal{X}$, l'espace algébrique $U$ possède la propriété $\mathcal{P}$.
\end{enumerate}
Si $\mathcal{X}$ est un schéma $U$, cela équivaut à $\mathcal{P}(U)$.
Si $\mathcal{X}$ est un espace algébrique $X$, cela équivaut à ce que
$X$ possède la propriété $\mathcal{P}$.
\end{lemma}

\begin{proof}
Soit $U \to \mathcal{X}$ un morphisme surjectif et lisse, où $U$ est un espace algébrique.
Soit $V \to \mathcal{X}$ un morphisme lisse, où $V$ est un espace algébrique.
Choisissons des schémas $U'$ et $V'$ et des morphismes \'etales surjectifs
$U' \to U$ et $V' \to V$. Enfin, choisissons un schéma $W$ et un
morphisme \'etale surjectif $W \to V' \times_\mathcal{X} U'$.
Alors $W \to V'$ et $W \to U'$ sont des morphismes lisses de schémas,
car ce sont des composés de morphismes \'etales et lisses d'espaces algébriques, voir
Morphisms of Spaces, Lemmes \ref{spaces-morphisms-lemma-etale-smooth} et
\ref{spaces-morphisms-lemma-composition-smooth}.
De plus, $W \to V'$ est surjectif puisque $U' \to \mathcal{X}$ est surjectif.
Nous avons donc
$$
\mathcal{P}(U) \Leftrightarrow
\mathcal{P}(U') \Rightarrow
\mathcal{P}(W) \Rightarrow
\mathcal{P}(V') \Leftrightarrow \mathcal{P}(V)
$$
où les équivalences résultent de la définition de la propriété $\mathcal{P}$ pour les
espaces algébriques, et les deux implications résultent de
Descent, Définition \ref{descent-definition-property-local}.
Cela démontre (3) $\Rightarrow$ (4).

\medskip\noindent
Les implications (2) $\Rightarrow$ (1), (1) $\Rightarrow$ (3),
et (4) $\Rightarrow$ (2) sont immédiates.
\end{proof}

\begin{definition}
\label{definition-type-property}
Soit $\mathcal{X}$ un champ algébrique.
Soit $\mathcal{P}$ une propriété des schémas qui est
locale pour la topologie lisse.
Nous disons que $\mathcal{X}$ {\it possède la propriété $\mathcal{P}$}
si l'une quelconque des conditions équivalentes du
Lemme \ref{lemma-type-property}
est satisfaite.
\end{definition}

\begin{remark}
\label{remark-list-properties-local-smooth-topology}
Voici une liste de propriétés qui sont locales pour la topologie lisse
(gardons à l'esprit que les topologies fpqc, fppf et syntomique sont
plus fines que la topologie lisse) :
\begin{enumerate}
\item localement noethérien, voir
Descent, Lemme \ref{descent-lemma-Noetherian-local-fppf},
\item de Jacobson, voir
Descent, Lemme \ref{descent-lemma-Jacobson-local-fppf},
\item localement noethérien et $(S_k)$, voir
Descent, Lemme \ref{descent-lemma-Sk-local-syntomic},
\item Cohen-Macaulay, voir
Descent, Lemme \ref{descent-lemma-CM-local-syntomic},
\item réduit, voir
Descent, Lemme \ref{descent-lemma-reduced-local-smooth},
\item normal, voir
Descent, Lemme \ref{descent-lemma-normal-local-smooth},
\item localement noethérien et $(R_k)$, voir
Descent, Lemme \ref{descent-lemma-Rk-local-smooth},
\item régulier, voir
Descent, Lemme \ref{descent-lemma-regular-local-smooth},
\item de Nagata, voir
Descent, Lemme \ref{descent-lemma-Nagata-local-smooth}.
\end{enumerate}
\end{remark}

\noindent
Toute propriété des germes de schémas locale pour la topologie lisse donne naissance à une
propriété correspondante des champs algébriques. Remarquons qu'une propriété des germes locale pour la topologie lisse
est aussi locale pour la topologie \'etale. Ainsi, pour une propriété $P$ des germes de
schémas locale pour la topologie lisse, nous savons ce que signifie dire qu'un espace algébrique $X$ possède
la propriété $P$ en $x \in |X|$, voir
Properties of Spaces, Section \ref{section-types-properties}.

\begin{lemma}
\label{lemma-local-source-target-at-point}
Soit $\mathcal{X}$ un champ algébrique.
Soit $x \in |\mathcal{X}|$ un point de $\mathcal{X}$.
Soit $\mathcal{P}$ une propriété des germes de schémas locale pour la topologie lisse, voir
Descent, Définition \ref{descent-definition-local-at-point}.
Les assertions suivantes sont équivalentes :
\begin{enumerate}
\item pour tout morphisme lisse $U \to \mathcal{X}$, où $U$ est un schéma,
et tout $u \in U$ tel que $a(u) = x$, on a $\mathcal{P}(U, u)$,
\item il existe un morphisme lisse $U \to \mathcal{X}$, où $U$ est un schéma,
et un $u \in U$ tel que $a(u) = x$, pour lesquels on a $\mathcal{P}(U, u)$,
\item pour tout morphisme lisse $U \to \mathcal{X}$, où $U$ est un espace algébrique,
et tout $u \in |U|$ tel que $a(u) = x$, l'espace algébrique $U$ possède la propriété
$\mathcal{P}$ en $u$, et
\item il existe un morphisme lisse $U \to \mathcal{X}$, où $U$ est un
espace algébrique, et un $u \in |U|$ tel que $a(u) = x$, tels que l'espace algébrique
$U$ possède la propriété $\mathcal{P}$ en $u$.
\end{enumerate}
Si $\mathcal{X}$ est représentable, cela équivaut à
$\mathcal{P}(\mathcal{X}, x)$. Si $\mathcal{X}$ est un espace algébrique, alors
cela équivaut à ce que $\mathcal{X}$ possède la propriété $\mathcal{P}$ en $x$.
\end{lemma}

\begin{proof}
Soient $a : U \to \mathcal{X}$ et $u \in |U|$ comme en (3). Soit
$b : V \to \mathcal{X}$ un autre morphisme lisse, où $V$ est un espace
algébrique, et soit $v \in |V|$ tel que $b(v) = x$ également.
Choisissons un schéma $U'$, un morphisme \'etale $U' \to U$ et $u' \in U'$
d'image $u$. Choisissons un schéma $V'$, un morphisme \'etale $V' \to V$
et $v' \in V'$ d'image $v$. D'après le
Lemme \ref{lemma-points-cartesian},
il existe un point $\overline{w} \in |V' \times_\mathcal{X} U'|$
d'images $u'$ et $v'$. Choisissons un schéma $W$ et un morphisme \'etale surjectif
$W \to V' \times_\mathcal{X} U'$. Nous pouvons choisir un
$w \in |W|$ d'image $\overline{w}$ (voir
Properties of Spaces,
Lemme \ref{spaces-properties-lemma-characterize-surjective}).
Alors $W \to V'$ et $W \to U'$ sont des morphismes lisses de schémas,
car ce sont des composés de morphismes \'etales et lisses d'espaces algébriques, voir
Morphisms of Spaces, Lemmes \ref{spaces-morphisms-lemma-etale-smooth} et
\ref{spaces-morphisms-lemma-composition-smooth}.
Par conséquent,
$$
\mathcal{P}(U, u)
\Leftrightarrow
\mathcal{P}(U', u')
\Leftrightarrow
\mathcal{P}(W, w)
\Leftrightarrow
\mathcal{P}(V', v')
\Leftrightarrow
\mathcal{P}(V, v)
$$
Les deux équivalences extérieures résultent de
Properties of Spaces,
Définition \ref{spaces-properties-definition-property-at-point},
et les deux autres de la définition d'une propriété locale pour la topologie lisse
des germes de schémas. Cela démontre (4) $\Rightarrow$ (3).

\medskip\noindent
Les implications (1) $\Rightarrow$ (2), (2) $\Rightarrow$ (4),
et (3) $\Rightarrow$ (1) sont immédiates.
\end{proof}

\begin{definition}
\label{definition-property-at-point}
Soit $\mathcal{P}$ une propriété des germes de schémas qui est
locale pour la topologie lisse. Soit $\mathcal{X}$ un champ algébrique.
Soit $x \in |\mathcal{X}|$.
Nous disons que $\mathcal{X}$ {\it possède la propriété $\mathcal{P}$ en $x$}
si l'une quelconque des conditions équivalentes du
Lemme \ref{lemma-local-source-target-at-point}
est satisfaite.
\end{definition}










\section{Monomorphismes de champs algébriques}
\label{section-monomorphisms}

\noindent
Nous définissons comme suit un monomorphisme de champs algébriques.
Nous verrons dans le
Lemme \ref{lemma-monomorphism}
que cela est compatible avec la notion $2$-catégorique correspondante.

\begin{definition}
\label{definition-monomorphism}
Soit $f : \mathcal{X} \to \mathcal{Y}$ un morphisme de champs algébriques.
Nous disons que $f$ est un {\it monomorphisme}
s'il est représentable par des espaces algébriques et est un monomorphisme au sens de la
Section \ref{section-properties-morphisms}.
\end{definition}

\noindent
Commençons par quelques lemmes élémentaires.

\begin{lemma}
\label{lemma-base-change-monomorphism}
Soit $\mathcal{X} \to \mathcal{Y}$ un morphisme de champs algébriques.
Soit $\mathcal{Z} \to \mathcal{Y}$ un monomorphisme.
Alors $\mathcal{Z} \times_\mathcal{Y} \mathcal{X} \to \mathcal{X}$
est un monomorphisme.
\end{lemma}

\begin{proof}
Cela résulte de la discussion générale de la
Section \ref{section-properties-morphisms}.
\end{proof}

\begin{lemma}
\label{lemma-composition-monomorphism}
Les composés de monomorphismes de champs algébriques sont des monomorphismes.
\end{lemma}

\begin{proof}
Cela résulte de la discussion générale de la
Section \ref{section-properties-morphisms}
et de
Morphisms of Spaces,
Lemme \ref{spaces-morphisms-lemma-composition-monomorphism}.
\end{proof}

\begin{lemma}
\label{lemma-monomorphism}
Soit $f : \mathcal{X} \to \mathcal{Y}$ un morphisme de champs algébriques.
Les assertions suivantes sont équivalentes :
\begin{enumerate}
\item $f$ est un monomorphisme,
\item $f$ est pleinement fidèle,
\item la diagonale
$\Delta_f : \mathcal{X} \to \mathcal{X} \times_\mathcal{Y} \mathcal{X}$
est une équivalence, et
\item il existe un espace algébrique $W$ et un morphisme surjectif et plat
$W \to \mathcal{Y}$, localement de présentation finie, tels que
$V = \mathcal{X} \times_\mathcal{Y} W$ soit un espace algébrique et que le
morphisme $V \to W$ soit un monomorphisme d'espaces algébriques.
\end{enumerate}
\end{lemma}

\begin{proof}
L'équivalence de (1) et (4) résulte de la discussion générale de la
Section \ref{section-properties-morphisms},
et en particulier des
Lemmes \ref{lemma-check-representable-covering} et
\ref{lemma-check-property-covering}.

\medskip\noindent
L'équivalence de (2) et (3) est
Categories, Lemme \ref{categories-lemma-fully-faithful-diagonal-equivalence}.

\medskip\noindent
Supposons satisfaites les conditions équivalentes (2) et (3). Alors $f$ est représentable
par des espaces algébriques d'après
Algebraic Stacks,
Lemme \ref{algebraic-lemma-characterize-representable-by-algebraic-spaces}.
En outre, le lemme de $2$-Yoneda, combiné à la pleine fidélité,
implique que, pour tout schéma $T$, le foncteur
$$
\Mor(T, \mathcal{X})
\longrightarrow
\Mor(T, \mathcal{Y})
$$
est pleinement fidèle. Étant donné un morphisme $y : T \to \mathcal{Y}$, il existe donc,
à un unique $2$-isomorphisme près, au plus un morphisme $x : T \to \mathcal{X}$,
tel que $y \cong f \circ x$. En particulier, étant donné un morphisme de schémas
$h : T' \to T$, il existe au plus un relèvement
$\tilde h : T' \to T \times_\mathcal{Y} \mathcal{X}$ de $h$.
Ainsi, $T \times_\mathcal{Y} \mathcal{X} \to T$ est un monomorphisme
d'espaces algébriques, ce qui démontre (1).

\medskip\noindent
Enfin, supposons (1). Alors, pour tout schéma $T$ et tout morphisme
$y : T \to \mathcal{Y}$, le produit fibré $T \times_\mathcal{Y} \mathcal{X}$
est un espace algébrique, et $T \times_\mathcal{Y} \mathcal{X} \to T$
est un monomorphisme. Il existe donc exactement un couple, à isomorphisme unique près,
$(x, \alpha)$, où $x : T \to \mathcal{X}$ est un morphisme
et $\alpha : f \circ x \to y$ est un $2$-morphisme. En appliquant
le lemme de $2$-Yoneda, cela signifie exactement que $f$ est pleinement fidèle, c'est-à-dire
que (2) est satisfaite.
\end{proof}

\begin{lemma}
\label{lemma-monomorphism-injective-points}
\begin{slogan}
Les monomorphismes de champs sont injectifs sur les points.
\end{slogan}
Un monomorphisme de champs algébriques induit une application injective entre les
ensembles de points.
\end{lemma}

\begin{proof}
Soit $f : \mathcal{X} \to \mathcal{Y}$ un monomorphisme de champs algébriques.
Supposons que $x_i : \Spec(K_i) \to \mathcal{X}$ soient des morphismes tels que
$f \circ x_1$ et $f \circ x_2$ définissent le même élément de $|\mathcal{Y}|$.
En appliquant la définition, nous trouvons une extension commune $\Omega$ munie des
morphismes correspondants $c_i : \Spec(\Omega) \to \Spec(K_i)$ et d'un
$2$-isomorphisme $\beta : f \circ x_1 \circ c_1 \to f \circ x_1 \circ c_2$.
Comme $f$ est pleinement fidèle, voir
Lemme \ref{lemma-monomorphism},
nous pouvons relever $\beta$ en un isomorphisme
$\alpha : x_1 \circ c_1 \to x_1 \circ c_2$.
Ainsi, $x_1$ et $x_2$ définissent le même point de $|\mathcal{X}|$,
comme souhaité.
\end{proof}

\begin{lemma}
\label{lemma-monomorphism-diagonal}
Soient $\mathcal{X} \to \mathcal{X}' \to \mathcal{Y}$ des morphismes
de champs algébriques. Si $\mathcal{X} \to \mathcal{X}'$ est un monomorphisme,
alors le diagramme canonique
$$
\xymatrix{
\mathcal{X} \ar[r] \ar[d] &
\mathcal{X} \times_\mathcal{Y} \mathcal{X} \ar[d] \\
\mathcal{X}' \ar[r] &
\mathcal{X}' \times_\mathcal{Y} \mathcal{X}'
}
$$
est un carré cartésien.
\end{lemma}

\begin{proof}
Nous avons $\mathcal{X} = \mathcal{X} \times_{\mathcal{X}'} \mathcal{X}$
d'après le Lemme \ref{lemma-monomorphism}. Le résultat en découle alors en appliquant
Categories, Lemme \ref{categories-lemma-fibre-product-after-map}.
\end{proof}







\section{Immersions de champs algébriques}
\label{section-immersions}

\noindent
Les immersions de champs algébriques sont définies comme suit.

\begin{definition}
\label{definition-immersion}
Immersions.
\begin{enumerate}
\item Un morphisme de champs algébriques est appelé une {\it immersion ouverte}
s'il est représentable et est une immersion ouverte
au sens de la
Section \ref{section-properties-morphisms}.
\item Un morphisme de champs algébriques est appelé une {\it immersion fermée}
s'il est représentable et est une immersion fermée
au sens de la
Section \ref{section-properties-morphisms}.
\item Un morphisme de champs algébriques est appelé une {\it immersion}
s'il est représentable et est une immersion
au sens de la
Section \ref{section-properties-morphisms}.
\end{enumerate}
\end{definition}

\noindent
Ce n'est pas pour nous la manière la plus commode de concevoir les immersions.
Il nous est un peu plus commode de considérer une immersion
comme un morphisme de champs algébriques représentable par des espaces
algébriques et qui est une immersion au sens de la
Section \ref{section-properties-morphisms}.
Il en va de même pour les immersions fermées et ouvertes.
Comme cela équivaut clairement à la notion qui vient d'être définie, nous
utiliserons cette caractérisation sans autre mention.
Nous démontrons quelques lemmes simples sur cette notion.

\begin{lemma}
\label{lemma-base-change-immersion}
Soit $\mathcal{X} \to \mathcal{Y}$ un morphisme de champs algébriques.
Soit $\mathcal{Z} \to \mathcal{Y}$ une immersion
(fermée, resp.\ ouverte).
Alors $\mathcal{Z} \times_\mathcal{Y} \mathcal{X} \to \mathcal{X}$
est une immersion (fermée, resp.\ ouverte).
\end{lemma}

\begin{proof}
Cela résulte de la discussion générale de la
Section \ref{section-properties-morphisms}.
\end{proof}

\begin{lemma}
\label{lemma-composition-immersion}
Les composés d'immersions de champs algébriques sont des immersions.
Il en va de même des immersions fermées et des immersions ouvertes.
\end{lemma}

\begin{proof}
Cela résulte de la discussion générale de la
Section \ref{section-properties-morphisms}
et de
Spaces, Lemme \ref{spaces-lemma-composition-immersions}.
\end{proof}

\begin{lemma}
\label{lemma-check-immersion-covering}
Soit $f : \mathcal{X} \to \mathcal{Y}$ un morphisme de champs algébriques.
Soit $W$ un espace algébrique et soit $W \to \mathcal{Y}$ un morphisme surjectif
et plat, localement de présentation finie. Les assertions suivantes
sont équivalentes :
\begin{enumerate}
\item $f$ est une immersion (ouverte, resp.\ fermée), et
\item $V = W \times_\mathcal{Y} \mathcal{X}$ est un espace algébrique, et
$V \to W$ est une immersion (ouverte, resp.\ fermée).
\end{enumerate}
\end{lemma}

\begin{proof}
Cela résulte de la discussion générale de la
Section \ref{section-properties-morphisms}
et en particulier des
Lemmes \ref{lemma-check-representable-covering} et
\ref{lemma-check-property-covering}.
\end{proof}

\begin{lemma}
\label{lemma-immersion-monomorphism}
Une immersion est un monomorphisme.
\end{lemma}

\begin{proof}
Voir
Morphisms of Spaces,
Lemme \ref{spaces-morphisms-lemma-immersions-monomorphisms}.
\end{proof}

\begin{lemma}
\label{lemma-subspace-induced-topology}
Si $f : \mathcal{X} \to \mathcal{Y}$ est une immersion, alors
$|f| : |\mathcal{X}| \to |\mathcal{Y}|$ est un homéomorphisme
sur une partie localement fermée. Si $f$ est une immersion fermée, resp.\ ouverte,
alors $|f|$ est fermée, resp.\ ouverte.
\end{lemma}

\begin{proof}
Omis.
\end{proof}

\noindent
Les deux lemmes suivants expliquent comment considérer les
immersions à l'aide de présentations.

\begin{lemma}
\label{lemma-immersion-into-presentation}
Soit $(U, R, s, t, c)$ un groupoïde lisse en espaces algébriques.
Soit $i : \mathcal{Z} \to [U/R]$ une immersion.
Il existe alors un sous-espace localement fermé $R$-invariant
$Z \subset U$ et une présentation $[Z/R_Z] \to \mathcal{Z}$,
où $R_Z$ est la restriction de $R$ à $Z$, tels que
$$
\xymatrix{
[Z/R_Z] \ar[dr] \ar[rr] & & \mathcal{Z} \ar[ld]^i \\
& [U/R]
}
$$
soit $2$-commutatif. Si $i$ est une immersion fermée (resp.\ ouverte),
alors $Z$ est un sous-espace fermé (resp.\ ouvert) de $U$.
\end{lemma}

\begin{proof}
D'après le
Lemme \ref{lemma-representable-in-terms-presentations},
nous obtenons un diagramme commutatif
$$
\xymatrix{
[U'/R'] \ar[dr] \ar[rr] & & \mathcal{Z} \ar[ld] \\
& [U/R]
}
$$
où $U' = \mathcal{Z} \times_{[U/R]} U$ et
$R' = \mathcal{Z} \times_{[U/R]} R$.
Puisque $\mathcal{Z} \to [U/R]$ est une immersion, nous voyons que
$U' \to U$ est une immersion d'espaces algébriques. Soit $Z \subset U$
le sous-espace localement fermé tel que $U' \to U$ se factorise par
$Z$ et induise un isomorphisme $U' \to Z$.
Il ressort clairement de la construction de $R'$ que
$R' = U' \times_{U, t} R = R \times_{s, U} U'$.
Cela implique que $Z \cong U'$ est $R$-invariant et que l'image de
$R' \to R$ identifie $R'$ à la restriction
$R_Z = s^{-1}(Z) = t^{-1}(Z)$ de $R$ à $Z$. Le lemme en résulte.
\end{proof}

\begin{lemma}
\label{lemma-immersion-presentation}
Soit $(U, R, s, t, c)$ un groupoïde lisse en espaces algébriques.
Soit $\mathcal{X} = [U/R]$ le champ algébrique associé, voir
Champs algébriques,
Théorème \ref{algebraic-theorem-smooth-groupoid-gives-algebraic-stack}.
Soit $Z \subset U$ un sous-espace localement fermé $R$-invariant. Alors
$$
[Z/R_Z] \longrightarrow [U/R]
$$
est une immersion de champs algébriques, où $R_Z$ est la restriction
de $R$ à $Z$. Si $Z \subset U$ est ouvert (resp.\ fermé), alors le morphisme
est une immersion ouverte (resp.\ fermée) de champs algébriques.
\end{lemma}

\begin{proof}
Rappelons que, d'après
Groupoïdes en espaces,
Définition \ref{spaces-groupoids-definition-invariant-open}
(voir aussi la discussion qui suit la définition),
on a $R_Z = s^{-1}(Z) = t^{-1}(Z)$ comme sous-espaces localement fermés
de $R$. Les deux morphismes $R_Z \to Z$ sont donc lisses comme changements de base
de $s$ et $t$. Par conséquent,
$(Z, R_Z, s|_{R_Z}, t|_{R_Z}, c|_{R_Z \times_{s, Z, t} R_Z})$ est
un groupoïde lisse en espaces algébriques, et l'on voit que
$[Z/R_Z]$ est un champ algébrique, voir
Champs algébriques,
Théorème \ref{algebraic-theorem-smooth-groupoid-gives-algebraic-stack}.
Les hypothèses de
Groupoïdes en espaces,
Lemme \ref{spaces-groupoids-lemma-criterion-fibre-product}
sont toutes satisfaites; il en résulte que l'on a un carré $2$-cartésien
$$
\xymatrix{
Z \ar[d] \ar[r] & [Z/R_Z] \ar[d] \\
U \ar[r] & [U/R]
}
$$
Il résulte de ceci et du
Lemme \ref{lemma-check-representable-covering}
que $[Z/R_Z] \to [U/R]$ est représentable par des espaces algébriques;
alors, il résulte du
Lemme \ref{lemma-check-property-covering}
que la flèche verticale de droite est une immersion (resp.\ une immersion fermée,
resp.\ une immersion ouverte) si et seulement si la flèche verticale de gauche l'est.
\end{proof}

\noindent
Nous pouvons définir les sous-champs ouverts, fermés et localement fermés comme suit.

\begin{definition}
\label{definition-substacks}
Soit $\mathcal{X}$ un champ algébrique.
\begin{enumerate}
\item Un {\it sous-champ ouvert} de $\mathcal{X}$ est une sous-catégorie strictement pleine
$\mathcal{X}' \subset \mathcal{X}$ telle que $\mathcal{X}'$ soit un champ
algébrique et que $\mathcal{X}' \to \mathcal{X}$ soit une immersion ouverte.
\item Un {\it sous-champ fermé} de $\mathcal{X}$ est une sous-catégorie strictement pleine
$\mathcal{X}' \subset \mathcal{X}$ telle que $\mathcal{X}'$ soit un champ
algébrique et que $\mathcal{X}' \to \mathcal{X}$ soit une immersion fermée.
\item Un {\it sous-champ localement fermé} de $\mathcal{X}$ est une sous-catégorie
strictement pleine $\mathcal{X}' \subset \mathcal{X}$ telle que $\mathcal{X}'$
soit un champ algébrique et que $\mathcal{X}' \to \mathcal{X}$ soit une immersion.
\end{enumerate}
\end{definition}

\noindent
Cette définition doit être employée avec précaution. En effet, si
$f : \mathcal{X} \to \mathcal{Y}$ est une équivalence de champs algébriques
et si $\mathcal{X}' \subset \mathcal{X}$ est un sous-champ ouvert, la sous-catégorie
$f(\mathcal{X}')$ n'est pas nécessairement un sous-champ ouvert de $\mathcal{Y}$.
Le problème vient de ce qu'elle peut ne pas être une sous-catégorie {\it strictement}
pleine; mais c'est aussi le seul problème.
Voici un énoncé formel.

\begin{lemma}
\label{lemma-substack-image}
Pour toute immersion $i : \mathcal{Z} \to \mathcal{X}$, il existe un
unique sous-champ localement fermé $\mathcal{X}' \subset \mathcal{X}$
tel que $i$ se factorise en la composée d'une
équivalence $i' : \mathcal{Z} \to \mathcal{X}'$
et du morphisme d'inclusion $\mathcal{X}' \to \mathcal{X}$.
Si $i$ est une immersion fermée (resp.\ ouverte), alors $\mathcal{X}'$
est un sous-champ fermé (resp.\ ouvert) de $\mathcal{X}$.
\end{lemma}

\begin{proof}
Omis.
\end{proof}

\begin{lemma}
\label{lemma-substacks-presentation}
Soit $[U/R] \to \mathcal{X}$ une présentation d'un champ algébrique.
Il existe une bijection canonique
$$
\text{sous-champs localement fermés }\mathcal{Z}\text{ de }\mathcal{X}
\longrightarrow
R\text{-invariants parmi les sous-espaces localement fermés }Z\text{ de }U
$$
qui envoie $\mathcal{Z}$ sur $U \times_\mathcal{X} \mathcal{Z}$.
De plus, un morphisme de champs algébriques $f : \mathcal{Y} \to \mathcal{X}$
se factorise par $\mathcal{Z}$ si et seulement si
$\mathcal{Y} \times_\mathcal{X} U \to U$ se factorise par $Z$.
Il en va de même des sous-champs fermés et des sous-champs ouverts.
\end{lemma}

\begin{proof}
D'après les Lemmes \ref{lemma-immersion-into-presentation} et
\ref{lemma-immersion-presentation},
l'application est une bijection.
Si $\mathcal{Y} \to \mathcal{X}$ se factorise par $\mathcal{Z}$,
le changement de base $\mathcal{Y} \times_\mathcal{X} U \to U$
se factorise bien sûr par $Z$. Réciproquement, supposons que $\mathcal{Y} \to \mathcal{X}$
soit un morphisme tel que $\mathcal{Y} \times_\mathcal{X} U \to U$
se factorise par $Z$. Nous allons montrer que, pour tout schéma $T$ et tout
morphisme $T \to \mathcal{Y}$,
donné par un objet $y$ de la catégorie fibre de $\mathcal{Y}$ au-dessus de $T$,
l'objet $y$ appartient en fait à la catégorie fibre de $\mathcal{Z}$ au-dessus de $T$.
En effet, le produit fibré $T \times_\mathcal{X} U$ est un espace
algébrique et $T \times_\mathcal{X} U \to T$ est un morphisme lisse surjectif.
Il existe donc un recouvrement fppf $\{T_i \to T\}$ tel que
$T_i \to T$ se factorise par $T \times_\mathcal{X} U \to T$ pour tout $i$.
Alors $T_i \to \mathcal{X}$ se factorise par $\mathcal{Y} \times_\mathcal{X} U$,
et donc par $Z \subset U$. Ainsi, $y|_{T_i}$ est un objet de
$\mathcal{Z}$ (car $Z$ est le produit fibré de $U$ et de
$\mathcal{Z}$ au-dessus de $\mathcal{X}$).
Puisque $\mathcal{Z}$ est un sous-champ strictement plein, nous concluons
que $y$ est un objet de $\mathcal{Z}$, comme voulu.
\end{proof}

\begin{lemma}
\label{lemma-open-substacks}
Soit $\mathcal{X}$ un champ algébrique. La règle
$\mathcal{U} \mapsto |\mathcal{U}|$ définit une bijection qui préserve les inclusions
entre les sous-champs ouverts de $\mathcal{X}$ et les parties ouvertes
de $|\mathcal{X}|$.
\end{lemma}

\begin{proof}
Choisissons une présentation $[U/R] \to \mathcal{X}$, voir
Champs algébriques, Lemme \ref{algebraic-lemma-stack-presentation}.
D'après le
Lemme \ref{lemma-substacks-presentation},
les sous-champs ouverts correspondent aux sous-schémas ouverts $R$-invariants
de $U$. D'autre part, les
Lemmes \ref{lemma-points-presentation} et \ref{lemma-topology-points}
assurent que ceux-ci correspondent bijectivement aux parties ouvertes de $|\mathcal{X}|$.
\end{proof}

\begin{lemma}
\label{lemma-open-image-substack}
Soit $\mathcal X$ un champ algébrique. Soit $U$ un espace algébrique et
$U \to \mathcal X$ un morphisme lisse surjectif. Pour une immersion ouverte
$V \hookrightarrow U$, il existe un champ algébrique $\mathcal Y$, une
immersion ouverte $\mathcal Y \to \mathcal X$ et un morphisme lisse
surjectif $V \to \mathcal Y$.
\end{lemma}

\begin{proof}
Nous définissons une catégorie fibrée en groupoïdes $\mathcal Y$ en prenant pour catégorie
fibre $\mathcal{Y}_T$ au-dessus d'un objet $T$ de $(\Sch/S)_{fppf}$
la sous-catégorie pleine de $\mathcal{X}_T$ formée de tous les
$y \in \Ob(\mathcal{X}_T)$ tels que le morphisme de projection
$V \times_{\mathcal X, y} T \to T$ soit surjectif. Or, pour tout morphisme
$x : T \to \mathcal X$, le produit $2$-fibré
$T \times_{x, \mathcal X} \mathcal Y$ a pour catégorie fibre au-dessus de $T'$ celle formée
des triplets $(f : T' \to T, y \in \mathcal{X}_{T'}, f^*x \simeq y)$ tels que
$V \times_{\mathcal X, y} T' \to T'$ soit surjectif.
Remarquons que $T \times_{x, \mathcal X} \mathcal Y$ est fibré en sétoïdes puisque
$\mathcal Y \to \mathcal X$ est fidèle (voir
Stacks, Lemme \ref{stacks-lemma-2-fibre-product-gives-stack-in-setoids}).
L'isomorphisme $f^*x \simeq y$ donne alors le diagramme
$$
\xymatrix{
V \times_{\mathcal X, y} T' \ar[d] \ar[r] &
V \times_{\mathcal X, x} T \ar[r] \ar[d] &
V \ar[d] \\
T' \ar[r]^f &
T \ar[r]^x &
\mathcal X
}
$$
où les deux carrés sont cartésiens. Le morphisme
$V \times_{\mathcal X, x} T \to T$ est lisse par changement de base, et donc ouvert.
Soit $T_0 \subset T$ son image. Les carrés cartésiens montrent que
$V \times_{\mathcal X, y} T' \to T'$ est surjectif si et seulement si $f$ se factorise
par $T_0$. Par conséquent, $T \times_{x, \mathcal X} \mathcal Y$ est représentable par
$T_0$, de sorte que l'inclusion $\mathcal Y \to \mathcal X$ est une immersion ouverte.
D'après
Champs algébriques, Lemme \ref{algebraic-lemma-open-fibred-category-is-algebraic},
nous concluons que $\mathcal{Y}$ est un champ algébrique.
Enfin, si nous désignons le morphisme $V \to \mathcal X$ par $g$, alors
$V \times_{\mathcal X} V \to V$ est surjectif (la diagonale fournit une
section). Ainsi, $g$ appartient à l'image de $\mathcal{Y}_V \to \mathcal{X}_V$, c'est-à-dire
que l'on obtient un morphisme $g' : V \to \mathcal{Y}$ s'inscrivant dans le
diagramme commutatif
$$
\xymatrix{
V \ar[r] \ar[d]^{g'} & U \ar[d] \\
\mathcal{Y} \ar[r] & \mathcal{X}
}
$$
Puisque $V \times_{g, \mathcal X} \mathcal Y \to V$ est un
monomorphisme, c'est en fait un isomorphisme car $(1, g')$ en définit une section.
Par conséquent, $g' : V \to \mathcal Y$ est un morphisme lisse, puisqu'il est le
changement de base du morphisme lisse $g : V \to \mathcal{X}$.
Il est surjectif par notre construction de $\mathcal{Y}$, ce qui achève
la démonstration du lemme.
\end{proof}

\begin{lemma}
\label{lemma-union-open-substacks}
Soit $\mathcal X$ un champ algébrique et soient $\mathcal{X}_i \subset \mathcal X$
une famille de sous-champs ouverts indexée par $i \in I$. Il existe alors un
sous-champ ouvert, que nous notons
$\bigcup_{i\in I} \mathcal{X}_i \subset \mathcal X$, tel que
les $\mathcal{X}_i$ soient des sous-champs ouverts qui le recouvrent.
\end{lemma}

\begin{proof}
Nous définissons une sous-catégorie fibrée
$\mathcal{X}' = \bigcup_{i \in I} \mathcal{X}_i$
en prenant pour catégorie fibre au-dessus d'un objet $T$ de $(\Sch/S)_{fppf}$
la sous-catégorie pleine de $\mathcal{X}_T$ formée de tous les
$x \in \Ob(\mathcal{X}_T)$ tels que le morphisme
$\coprod_{i \in I} (\mathcal{X}_i \times_{\mathcal X} T) \to T$
soit surjectif. Soit $x_i \in \Ob((\mathcal{X}_i)_T)$.
Alors $(x_i, 1)$ définit une section de
$\mathcal{X}_i \times_{\mathcal X} T \to T$, et donc un isomorphisme. Ainsi,
$\mathcal{X}_i \subset \mathcal{X}'$ est une sous-catégorie pleine.
Soit maintenant $x \in \Ob(\mathcal{X}_T)$. Alors
$\mathcal{X}_i \times_{\mathcal X} T$ est représentable
par un sous-schéma ouvert $T_i \subset T$. Le produit $2$-fibré
$\mathcal{X}' \times_{\mathcal X} T$ a pour fibre au-dessus de $T'$ les
$(y \in \mathcal{X}_{T'}, f : T' \to T, f^*x \simeq y)$ tels que
$\coprod (\mathcal{X}_i \times_{\mathcal X, y} T') \to T'$ soit surjectif.
L'isomorphisme $f^*x \simeq y$ induit un isomorphisme
$\mathcal{X}_i \times_{\mathcal X, y} T' \simeq T_i \times_T T'$.
Alors les $T_i \times_T T'$ recouvrent $T'$ si et seulement si $f$ se factorise par
$\bigcup T_i$. Nous avons donc un diagramme
$$
\xymatrix{
T_i \ar[r] \ar[d] &
\bigcup T_i \ar[r] \ar[d] &
T \ar[d] \\
\mathcal{X}_i \ar[r] &
\mathcal{X}' \ar[r] &
\mathcal{X}
}
$$
dont les deux carrés sont cartésiens. D'après
Champs algébriques, Lemme \ref{algebraic-lemma-open-fibred-category-is-algebraic},
nous concluons que $\mathcal{X'} \subset \mathcal{X}$ est algébrique et constitue un
sous-champ ouvert. Il ressort aussi des carrés cartésiens ci-dessus que le
morphisme $\coprod_{i \in I} \mathcal{X}_i \to \mathcal{X}'$ est surjectif, ce
qui achève la démonstration du lemme.
\end{proof}

\begin{lemma}
\label{lemma-quasi-compact-finite-subcover}
Soit $\mathcal X$ un champ algébrique et soit $\mathcal X' \subset \mathcal X$
un sous-champ ouvert quasi-compact. Supposons donnée une famille de sous-
champs ouverts $\mathcal{X}_i \subset \mathcal X$ indexée par $i \in I$ telle
que $\mathcal{X}' \subset \bigcup_{i \in I} \mathcal{X}_i$,
où la réunion est définie comme au Lemme \ref{lemma-union-open-substacks}.
Il existe alors une partie finie $I' \subset I$ telle que
$\mathcal{X}' \subset \bigcup_{i \in I'} \mathcal{X}_i$.
\end{lemma}

\begin{proof}
Puisque $\mathcal X$ est algébrique, il existe un schéma $U$ muni d'un morphisme
lisse surjectif $U \to \mathcal X$. Soit $U_i \subset U$ le sous-schéma ouvert
représentant $\mathcal{X}_i \times_{\mathcal X} U$, et $U' \subset U$ le
sous-schéma ouvert représentant $\mathcal{X}' \times_{\mathcal X} U$. Par
hypothèse, $U'\subset \bigcup_{i\in I} U_i$. D'après la démonstration du
Lemme \ref{lemma-quasi-compact-stack},
il existe un ouvert quasi-compact $V \subset U'$ tel que $V \to \mathcal{X}'$
soit un morphisme lisse surjectif. Il existe donc une partie finie
$I' \subset I$ telle que $V \subset \bigcup_{i \in I'}
U_i$. Nous affirmons que $\mathcal{X}' \subset \bigcup_{i \in I'} \mathcal{X}_i$.
Prenons $x \in \Ob(\mathcal{X}'_T)$ pour
$T \in \Ob((\Sch/S)_{fppf})$.
Comme $\mathcal{X}' \to \mathcal{X}$ est un monomorphisme, nous avons les carrés
cartésiens
$$
\xymatrix{
V \times_\mathcal{X} T \ar[r] \ar[d] &
T \ar[d]^x \ar@{=}[r] &
T \ar[d]^x \\
V \ar[r] &
\mathcal{X}' \ar[r] &
\mathcal X
}
$$
Par changement de base, $V \times_{\mathcal X} T \to T$ est surjectif. Par conséquent,
$\bigcup_{i \in I'} U_i \times_{\mathcal X} T \to T$ est lui aussi surjectif.
Soit $T_i \subset T$ le sous-schéma ouvert représentant
$\mathcal{X}_i \times_{\mathcal X} T$.
Un argument formel donne le carré cartésien
$$
\xymatrix{
U_i \times_{\mathcal{X}_i} T_i \ar[r] \ar[d] &
U \times_{\mathcal X} T \ar[d] \\
T_i \ar[r] & T
}
$$
où les flèches verticales sont surjectives par changement de base. Puisque
$U_i \times_{\mathcal{X}_i} T_i \simeq U_i \times_{\mathcal X} T$,
nous obtenons $\bigcup_{i \in I'} T_i = T$. Ainsi,
$x$ est un objet de $(\bigcup_{i\in I'} \mathcal{X}_i)_T$ par
définition de la réunion. Observons que l'inclusion
$\mathcal{X}' \subset \bigcup_{i \in I'} \mathcal{X}_i$
est automatiquement un sous-champ ouvert.
\end{proof}

\begin{lemma}
\label{lemma-zariski-open-cover-stack-is-space}
Soit $\mathcal X$ un champ algébrique.
Soit $\mathcal{X}_i$, $i \in I$, une famille de sous-champs ouverts de $\mathcal{X}$.
Supposons que
\begin{enumerate}
\item $\mathcal{X} = \bigcup_{i \in I} \mathcal{X}_i$, et que
\item chaque $\mathcal{X}_i$ soit un espace algébrique.
\end{enumerate}
Alors $\mathcal{X}$ est un espace algébrique.
\end{lemma}

\begin{proof}
Appliquons
Champs, Lemme \ref{stacks-lemma-stack-in-setoids-descent}
au morphisme $\coprod_{i \in I} \mathcal{X}_i \to \mathcal{X}$
et au morphisme $\text{id} : \mathcal{X} \to \mathcal{X}$ pour
voir que $\mathcal{X}$ est un champ en sétoïdes.
Ainsi, $\mathcal{X}$ est un espace algébrique, voir
Champs algébriques,
Proposition \ref{algebraic-proposition-algebraic-stack-no-automorphisms}.
\end{proof}

\begin{lemma}
\label{lemma-zariski-open-cover-stack-is-scheme}
Soit $\mathcal X$ un champ algébrique.
Soit $\mathcal{X}_i$, $i \in I$, une famille de sous-champs ouverts de $\mathcal{X}$.
Supposons que
\begin{enumerate}
\item $\mathcal{X} = \bigcup_{i \in I} \mathcal{X}_i$, et que
\item chaque $\mathcal{X}_i$ soit un schéma.
\end{enumerate}
Alors $\mathcal{X}$ est un schéma.
\end{lemma}

\begin{proof}
D'après le
Lemme \ref{lemma-zariski-open-cover-stack-is-space},
on voit que $\mathcal{X}$ est un espace algébrique. Comme tout espace
algébrique possède un plus grand sous-espace ouvert qui soit un schéma, voir
Propriétés des espaces, Lemme \ref{spaces-properties-lemma-subscheme},
on voit que $\mathcal{X}$ est un schéma.
\end{proof}


\noindent
Le lemme suivant est l'analogue de
Compléments sur les groupoïdes, Lemme \ref{more-groupoids-lemma-local-source}.

\begin{lemma}
\label{lemma-local-source}
Soient $\mathcal{P}, \mathcal{Q}, \mathcal{R}$ des propriétés de morphismes
d'espaces algébriques. Supposons que
\begin{enumerate}
\item $\mathcal{P}, \mathcal{Q}, \mathcal{R}$ soient locales pour la topologie fppf sur le but
et stables par tout changement de base,
\item $\text{lisse} \Rightarrow \mathcal{R}$,
\item pour tout morphisme $f : X \to Y$ qui possède $\mathcal{Q}$, il existe un
plus grand sous-espace ouvert $W(\mathcal{P}, f) \subset X$ tel que
$f|_{W(\mathcal{P}, f)}$ possède $\mathcal{P}$, et
\item pour tout morphisme $f : X \to Y$ qui possède $\mathcal{Q}$,
et tout morphisme $Y' \to Y$ qui possède $\mathcal{R}$, on ait
$Y' \times_Y W(\mathcal{P}, f) = W(\mathcal{P}, f')$, où
$f' : X_{Y'} \to Y'$ est le changement de base de $f$.
\end{enumerate}
Soit $f : \mathcal{X} \to \mathcal{Y}$ un morphisme de champs algébriques
représentable par des espaces algébriques. Supposons que $f$ possède $\mathcal{Q}$. Alors
\begin{enumerate}
\item[(A)] il existe un plus grand sous-champ ouvert
$\mathcal{X}' \subset \mathcal{X}$ tel que $f|_{\mathcal{X}'}$ possède
$\mathcal{P}$, et
\item[(B)] si $\mathcal{Z} \to \mathcal{Y}$ est un morphisme de champs
algébriques représentable par des espaces algébriques qui possède $\mathcal{R}$,
alors $\mathcal{Z} \times_\mathcal{Y} \mathcal{X}'$ est le plus grand sous-champ
ouvert de $\mathcal{Z} \times_\mathcal{Y} \mathcal{X}$ au-dessus duquel
le changement de base $\text{id}_\mathcal{Z} \times f$ possède la propriété $\mathcal{P}$.
\end{enumerate}
\end{lemma}

\begin{proof}
Choisissons un schéma $V$ et un morphisme lisse surjectif $V \to \mathcal{Y}$.
Posons $U = V \times_\mathcal{Y} \mathcal{X}$ et soit $f' : U \to V$ le
changement de base de $f$. Le morphisme d'espaces algébriques $f' : U \to V$ possède
la propriété $\mathcal{Q}$. L'hypothèse (3) fournit donc l'ouvert $W(\mathcal{P}, f') \subset U$.
Remarquons que
$U \times_\mathcal{X} U =
(V \times_\mathcal{Y} V) \times_\mathcal{Y} \mathcal{X}$,
si bien que le morphisme $f'' : U \times_\mathcal{X} U \to V \times_\mathcal{Y} V$
est le changement de base de $f$ par l'une ou l'autre projection
$V \times_\mathcal{Y} V \to V$. Par notre choix de $V$, ces projections
sont lisses, et possèdent donc la propriété $\mathcal{R}$ d'après (2). Ainsi, (4) montre
que les images inverses de $W(\mathcal{P}, f')$ par les deux projections
$\text{pr}_i : U \times_\mathcal{X} U \to U$ coïncident. Autrement dit,
$W(\mathcal{P}, f')$ est un sous-espace $R$-invariant de $U$ (où
$R = U \times_\mathcal{X} U$). Soit $\mathcal{X}'$ le sous-champ ouvert de
$\mathcal{X}$ qui correspond à $W(\mathcal{P}, f)$ par le
Lemme \ref{lemma-immersion-into-presentation}.
Par construction, $W(\mathcal{P}, f') = \mathcal{X}' \times_\mathcal{Y} V$,
donc $f|_{\mathcal{X}'}$ possède la propriété $\mathcal{P}$ d'après le
Lemme \ref{lemma-check-property-covering}.
De plus, $\mathcal{X}'$ est le plus grand sous-champ ouvert
tel que $f|_{\mathcal{X}'}$ possède $\mathcal{P}$, puisque la même maximalité vaut
pour $W(\mathcal{P}, f)$. Cela prouve (A).

\medskip\noindent
Enfin, si $\mathcal{Z} \to \mathcal{Y}$ est un
morphisme de champs algébriques représentable par des espaces algébriques qui possède
$\mathcal{R}$, posons $T = V \times_\mathcal{Y} \mathcal{Z}$; on voit alors que
$T \to V$ est un morphisme d'espaces algébriques possédant
la propriété $\mathcal{R}$. Soit $f'_T : T \times_V U \to T$ le changement de base
de $f'$. De nouveau, (4) montre que $W(\mathcal{P}, f'_T)$
est l'image inverse de $W(\mathcal{P}, f)$ dans $T \times_V U$. Cela implique
(B); nous omettons certains détails.
\end{proof}

\begin{remark}
\label{remark-local-source-warning}
Avertissement :
Le Lemme \ref{lemma-local-source}
doit être utilisé avec précaution. Par exemple, il s'applique à
$\mathcal{P}=$``plat'', $\mathcal{Q}=$``vide'',
et $\mathcal{R}=$``plat et localement de présentation finie''. Mais, étant donné un
morphisme d'espaces algébriques $f : X \to Y$, le plus grand sous-
espace ouvert $W \subset X$ tel que $f|_W$ soit plat n'est {\it pas} l'ensemble des points
où $f$ est plat !
\end{remark}

\begin{remark}
\label{remark-local-source-apply}
Malgré l'avertissement de la
Remarque \ref{remark-local-source-warning},
il existe des cas où le
Lemme \ref{lemma-local-source}
peut être utilisé sans ambiguïté.
Nous en donnons une liste. Dans chaque cas, nous omettons la vérification des
hypothèses (1) et (2), et donnons des références qui impliquent
(3) et (4). Voici la liste :
\begin{enumerate}
\item
\label{item-rel-dim-leq-d}
$\mathcal{Q} = $``localement de type fini'', $\mathcal{R} = \emptyset$,
et $\mathcal{P} =$``dimension relative $\leq d$''.
Voir
Morphismes d'espaces,
Définition \ref{spaces-morphisms-definition-relative-dimension}
et
Morphismes d'espaces, Lemmes
\ref{spaces-morphisms-lemma-openness-bounded-dimension-fibres} et
\ref{spaces-morphisms-lemma-dimension-fibre-after-base-change}.
\item
\label{item-loc-quasi-finite}
$\mathcal{Q} =$``localement de type fini'', $\mathcal{R} = \emptyset$,
et $\mathcal{P} =$``localement quasi-fini''.
Il s'agit du cas $d = 0$ du point précédent, voir
Morphismes d'espaces, Lemme
\ref{spaces-morphisms-lemma-locally-quasi-finite-rel-dimension-0}.
D'autre part, les propriétés (3) et (4) sont explicitées dans
Morphismes d'espaces, Lemme
\ref{spaces-morphisms-lemma-locally-finite-type-quasi-finite-part}.
\item
\label{item-unramified}
$\mathcal{Q} = $``localement de type fini'', $\mathcal{R} = \emptyset$,
et $\mathcal{P} =$``non ramifié''. C'est le
Morphismes d'espaces, Lemme \ref{spaces-morphisms-lemma-where-unramified}.
\item
\label{item-flat}
$\mathcal{Q} =$``localement de présentation finie'',
$\mathcal{R} =$``plat et localement de présentation finie'', et
$\mathcal{P} =$``plat''. Voir
Compléments sur les morphismes d'espaces, Théorème
\ref{spaces-more-morphisms-theorem-openness-flatness} et
Lemme \ref{spaces-more-morphisms-lemma-flat-locus-base-change}.
Remarquons qu'ici $W(\mathcal{P}, f)$ est toujours exactement l'ensemble des points
où le morphisme $f$ est plat, car nous ne considérons cet ouvert que
lorsque $f$ possède $\mathcal{Q}$ (voir loc. cit.).
\item
\label{item-etale}
$\mathcal{Q} =$``localement de présentation finie'',
$\mathcal{R} =$``plat et localement de présentation finie'', et
$\mathcal{P}=$``\'etale''. Cela résulte de la combinaison de
(\ref{item-unramified}) et (\ref{item-flat}), puisqu'un morphisme non ramifié
qui est plat et localement de présentation finie est \'etale, voir
Morphismes d'espaces,
Lemme \ref{spaces-morphisms-lemma-unramified-flat-lfp-etale}.
\item Ajouter ici d'autres cas selon les besoins (comparer avec la liste plus longue de
Compléments sur les groupoïdes, Remarque \ref{more-groupoids-remark-local-source-apply}).
\end{enumerate}
\end{remark}











\section{Champs algébriques réduits}
\label{section-reduced}

\noindent
Nous avons déjà défini les champs algébriques réduits dans la
Section \ref{section-types-properties}.

\begin{lemma}
\label{lemma-reduced-closed-substack}
Soit $\mathcal{X}$ un champ algébrique.
Soit $T \subset |\mathcal{X}|$ une partie fermée.
Il existe un unique sous-champ fermé $\mathcal{Z} \subset \mathcal{X}$
possédant les propriétés suivantes :
(a) on a $|\mathcal{Z}| = T$, et (b) $\mathcal{Z}$ est réduit.
\end{lemma}

\begin{proof}
Soit $U \to \mathcal{X}$ un morphisme lisse surjectif, où $U$ est un
espace algébrique. Posons $R = U \times_\mathcal{X} U$; on a alors une
présentation $[U/R] \to \mathcal{X}$, voir
Champs algébriques, Lemme \ref{algebraic-lemma-stack-presentation}.
Comme d'habitude, nous notons $s, t : R \to U$ les deux morphismes de projection lisses.
D'après le Lemme \ref{lemma-points-presentation},
on voit que $T$ correspond à une partie fermée $T' \subset |U|$ telle
que $|s|^{-1}(T') = |t|^{-1}(T')$.
Soit $Z \subset U$ l'espace algébrique obtenu en munissant
$T'$ de la structure réduite induite, voir
Propriétés des espaces,
Définition \ref{spaces-properties-definition-reduced-induced-space}.
Les produits fibrés
$Z \times_{U, t} R$ et $R \times_{s, U} Z$ sont des sous-espaces fermés
de $R$
(Espaces, Lemme \ref{spaces-lemma-base-change-immersions}).
Les projections $Z \times_{U, t} R \to Z$ et
$R \times_{s, U} Z \to Z$ sont lisses d'après
Morphismes d'espaces, Lemme \ref{spaces-morphisms-lemma-base-change-smooth}.
Ainsi, puisque $Z$ est réduit, il s'ensuit que
$Z \times_{U, t} R$ et $R \times_{s, U} Z$ sont réduits, voir
Remarque \ref{remark-list-properties-local-smooth-topology}.
Comme
$$
|Z \times_{U, t} R| = |t|^{-1}(T') = |s|^{-1}(T') = R \times_{s, U} Z
$$
l'unicité de
Propriétés des espaces,
Lemme \ref{spaces-properties-lemma-reduced-closed-subspace}
entraîne que $Z \times_{U, t} R = R \times_{s, U} Z$.
Ainsi, $Z$ est un sous-espace fermé $R$-invariant de $U$.
Par la correspondance du
Lemme \ref{lemma-substacks-presentation},
nous obtenons un sous-champ fermé $\mathcal{Z} \subset \mathcal{X}$
tel que $Z = \mathcal{Z} \times_\mathcal{X} U$. Alors
$[Z/R_Z] \to \mathcal{Z}$ est une présentation
(Lemme \ref{lemma-immersion-into-presentation}).
Enfin, $|\mathcal{Z}| = |Z|/|R_Z| = |T'|/\sim$ est la
partie fermée donnée $T$. Nous omettons la démonstration de l'unicité.
\end{proof}

\begin{lemma}
\label{lemma-reduced-stack-determined-by-points}
Soit $\mathcal{X}$ un champ algébrique.
Si $\mathcal{X}' \subset \mathcal{X}$
est un sous-champ fermé, si $\mathcal{X}$ est réduit et si
$|\mathcal{X}'| = |\mathcal{X}|$, alors $\mathcal{X}' = \mathcal{X}$.
\end{lemma}

\begin{proof}
Choisissons une présentation $[U/R] \to \mathcal{X}$ avec $U$ un schéma.
Comme $\mathcal{X}$ est réduit, $U$ est réduit (par définition
des champs algébriques réduits). D'après le
Lemme \ref{lemma-substacks-presentation},
$\mathcal{X}'$ correspond à un sous-schéma fermé $R$-invariant $Z \subset U$.
Mais $|Z| \subset |U|$ est l'image inverse de $|\mathcal{X}'|$, et
donc $|Z| = |U|$. Ainsi, $Z$ est un sous-schéma fermé de $U$ dont les ensembles
de points sous-jacents coïncident. D'après
Schémas, Lemme \ref{schemes-lemma-map-into-reduction},
l'application $\text{id}_U : U \to U$ se factorise par $Z \to U$, et donc
$Z = U$, c'est-à-dire $\mathcal{X}' = \mathcal{X}$.
\end{proof}

\begin{lemma}
\label{lemma-map-into-reduction}
Soient $\mathcal{X}$, $\mathcal{Y}$ des champs algébriques.
Soit $\mathcal{Z} \subset \mathcal{X}$ un sous-champ fermé.
Supposons $\mathcal{Y}$ réduit.
Un morphisme $f : \mathcal{Y} \to \mathcal{X}$ se factorise par
$\mathcal{Z}$ si et seulement si
$f(|\mathcal{Y}|) \subset |\mathcal{Z}|$.
\end{lemma}

\begin{proof}
Supposons $f(|\mathcal{Y}|) \subset |\mathcal{Z}|$. Considérons
$\mathcal{Y} \times_\mathcal{X} \mathcal{Z} \to \mathcal{Y}$.
Il existe une équivalence
$\mathcal{Y} \times_\mathcal{X} \mathcal{Z} \to \mathcal{Y}'$
où $\mathcal{Y}'$ est un sous-champ fermé de $\mathcal{Y}$, voir les
Lemmes \ref{lemma-base-change-immersion} et
\ref{lemma-substack-image}.
À l'aide des
Lemmes \ref{lemma-points-cartesian},
\ref{lemma-monomorphism-injective-points} et
\ref{lemma-immersion-monomorphism},
on voit que $|\mathcal{Y}'| = |\mathcal{Y}|$. Le lemme est donc ramené au
Lemme \ref{lemma-reduced-stack-determined-by-points}.
\end{proof}

\begin{definition}
\label{definition-reduced-induced-stack}
Soit $\mathcal{X}$ un champ algébrique.
Soit $Z \subset |\mathcal{X}|$ une partie fermée.
Une {\it structure de champ algébrique sur $Z$} est donnée par un sous-champ fermé
$\mathcal{Z}$ de $\mathcal{X}$ tel que $|\mathcal{Z}|$ soit égal à $Z$.
La {\it structure de champ algébrique réduite induite}
sur $Z$ est celle construite au
Lemme \ref{lemma-reduced-closed-substack}.
La {\it réduction $\mathcal{X}_{red}$ de $\mathcal{X}$}
est la structure de champ algébrique réduite induite sur $|\mathcal{X}|$.
\end{definition}


\noindent
En fait, cela permet de définir la structure de champ algébrique réduite
induite sur une partie localement fermée.

\begin{remark}
\label{remark-stack-structure-locally-closed-subset}
Soit $X$ un champ algébrique.
Soit $T \subset |\mathcal{X}|$ une partie localement fermée.
Soit $\partial T$ la frontière de $T$ dans
l'espace topologique $|\mathcal{X}|$. En formule,
$$
\partial T = \overline{T} \setminus T.
$$
Soit $\mathcal{U} \subset \mathcal{X}$ le sous-champ ouvert de $X$ tel que
$|\mathcal{U}| = |\mathcal{X}| \setminus \partial T$, voir
Lemme \ref{lemma-open-substacks}.
Soit $\mathcal{Z}$ le sous-champ fermé réduit de $\mathcal{U}$ tel que
$|\mathcal{Z}| = T$, obtenu en prenant la structure réduite induite
de sous-espace fermé, voir
Définition \ref{definition-reduced-induced-stack}.
Par construction, $\mathcal{Z} \to \mathcal{U}$ est une immersion fermée de
champs algébriques et $\mathcal{U} \to \mathcal{X}$ est une immersion ouverte;
ainsi, $\mathcal{Z} \to \mathcal{X}$ est une immersion de champs algébriques d'après le
Lemme \ref{lemma-composition-immersion}.
Remarquons que $\mathcal{Z}$ est un champ algébrique réduit et que
$|\mathcal{Z}| = T$ comme parties de $|X|$. On dit parfois que
$\mathcal{Z}$ est la {\it structure de sous-champ réduite induite} sur $T$.
\end{remark}













\section{Gerbes résiduelles}
\label{section-residual-gerbe}

\noindent
Dans le projet Stacks, nous voudrions définir la {\it gerbe résiduelle}
d'un champ algébrique $\mathcal{X}$ en un point $x \in |\mathcal{X}|$
comme un monomorphisme de champs algébriques
$m_x : \mathcal{Z}_x \to \mathcal{X}$ où $\mathcal{Z}_x$ est un champ
algébrique réduit ayant un unique point, envoyé par $m_x$ sur $x$.
Il s'avère que cette notion soulève de nombreuses difficultés; l'existence n'est pas
claire en général, pas plus que l'unicité.
Nous résolvons le problème d'unicité en imposant une condition légèrement plus forte
aux champs algébriques $\mathcal{Z}_x$.
Nous examinons cela plus en détail au moyen de quelques lemmes simples
sur les champs algébriques réduits ayant un unique point.

\begin{lemma}
\label{lemma-flat-cover-by-field}
Soit $\mathcal{Z}$ un champ algébrique. Soit $k$ un corps et soit
$\Spec(k) \to \mathcal{Z}$ un morphisme plat surjectif. Alors tout
morphisme $\Spec(k') \to \mathcal{Z}$, où $k'$ est un corps, est
plat et surjectif.
\end{lemma}

\begin{proof}
Considérons le carré fibré
$$
\xymatrix{
T \ar[d] \ar[r] & \Spec(k) \ar[d] \\
\Spec(k') \ar[r] & \mathcal{Z}
}
$$
Remarquons que $T \to \Spec(k')$ est plat et surjectif; ainsi, $T$
n'est pas vide. D'autre part, $T \to \Spec(k)$ est plat puisque
$k$ est un corps. Il s'ensuit que $T \to \mathcal{Z}$ est plat et surjectif.
D'après
Morphismes d'espaces, Lemme \ref{spaces-morphisms-lemma-flat-permanence}
(via la discussion de la
Section \ref{section-properties-morphisms}),
on en déduit que $\Spec(k') \to \mathcal{Z}$ est plat. Il est clairement
surjectif puisque, par hypothèse, $|\mathcal{Z}|$ est réduit à un point.
\end{proof}

\begin{lemma}
\label{lemma-unique-point}
Soit $\mathcal{Z}$ un champ algébrique. Les assertions suivantes sont équivalentes :
\begin{enumerate}
\item $\mathcal{Z}$ est réduit et $|\mathcal{Z}|$ est réduit à un point,
\item il existe un morphisme plat surjectif $\Spec(k) \to \mathcal{Z}$,
où $k$ est un corps, et
\item il existe un morphisme localement de type fini, surjectif et plat
$\Spec(k) \to \mathcal{Z}$, où $k$ est un corps.
\end{enumerate}
\end{lemma}

\begin{proof}
Supposons (1). Soit $W$ un schéma et
soit $W \to \mathcal{Z}$ un morphisme lisse surjectif. Alors $W$ est
un schéma réduit. Soit $\eta \in W$ un point générique d'une composante
irréductible de $W$. Comme $W$ est réduit, on a
$\mathcal{O}_{W, \eta} = \kappa(\eta)$. Il s'ensuit que le morphisme canonique
$\eta = \Spec(\kappa(\eta)) \to W$ est plat. On voit que le
composé $\eta \to \mathcal{Z}$ est plat (voir
Morphismes d'espaces, Lemme \ref{spaces-morphisms-lemma-composition-flat}).
Il est aussi surjectif puisque $|\mathcal{Z}|$ est réduit à un point. Autrement dit,
(2) est vérifiée.

\medskip\noindent
Supposons (2). Soit $W$ un schéma et
soit $W \to \mathcal{Z}$ un morphisme lisse surjectif. Choisissons un corps
$k$ et un morphisme plat surjectif $\Spec(k) \to \mathcal{Z}$.
Alors $W \times_\mathcal{Z} \Spec(k)$ est un espace algébrique lisse
sur $k$, donc régulier (voir
Espaces sur les corps, Lemme \ref{spaces-over-fields-lemma-smooth-regular})
et en particulier réduit. Comme $W \times_\mathcal{Z} \Spec(k) \to W$
est plat et surjectif, nous concluons que $W$ est réduit
(Descente sur les espaces, Lemme \ref{spaces-descent-lemma-descend-reduced}).
Autrement dit, (1) est vérifiée.

\medskip\noindent
Il est clair que (3) implique (2). Enfin, supposons (2). Prenons un schéma affine
non vide $W$ et un morphisme lisse $W \to \mathcal{Z}$. Prenons un point fermé
$w \in W$ et posons $k = \kappa(w)$. Le composé
$$
\Spec(k) \xrightarrow{w} W \longrightarrow \mathcal{Z}
$$
est localement de type fini d'après
Morphismes d'espaces, Lemmes
\ref{spaces-morphisms-lemma-composition-finite-type} et
\ref{spaces-morphisms-lemma-smooth-locally-finite-type}.
Il est aussi plat et surjectif d'après le
Lemme \ref{lemma-flat-cover-by-field}.
Ainsi, (3) est vérifiée.
\end{proof}

\noindent
Le lemme suivant distingue une classe légèrement meilleure de champs algébriques
réduits à un point que celle du lemme précédent.

\begin{lemma}
\label{lemma-unique-point-better}
Soit $\mathcal{Z}$ un champ algébrique. Les assertions suivantes sont équivalentes :
\begin{enumerate}
\item $\mathcal{Z}$ est réduit, localement noethérien, et $|\mathcal{Z}|$
est réduit à un point, et
\item il existe un morphisme localement de présentation finie, surjectif et plat
$\Spec(k) \to \mathcal{Z}$, où $k$ est un corps.
\end{enumerate}
\end{lemma}

\begin{proof}
Supposons (2). D'après le
Lemme \ref{lemma-unique-point},
on voit que $\mathcal{Z}$ est réduit et que $|\mathcal{Z}|$ est réduit à un point.
Soit $W$ un schéma et soit $W \to \mathcal{Z}$ un morphisme lisse
surjectif. Choisissons un corps $k$ et un morphisme localement de présentation finie,
surjectif et plat $\Spec(k) \to \mathcal{Z}$.
Alors $W \times_\mathcal{Z} \Spec(k)$ est un espace algébrique
lisse sur $k$, donc localement noethérien (voir
Morphismes d'espaces, Lemme
\ref{spaces-morphisms-lemma-locally-finite-type-locally-noetherian}).
Comme $W \times_\mathcal{Z} \Spec(k) \to W$
est plat, surjectif et localement de présentation finie, on voit
que $\{W \times_\mathcal{Z} \Spec(k) \to W\}$ est un recouvrement fppf,
et nous concluons que $W$ est localement noethérien
(Descente sur les espaces, Lemme
\ref{spaces-descent-lemma-descend-locally-Noetherian}).
Autrement dit, (1) est vérifiée.

\medskip\noindent
Supposons (1). Prenons un schéma affine non vide $W$ et un morphisme lisse
$W \to \mathcal{Z}$. Prenons un point fermé $w \in W$ et posons
$k = \kappa(w)$. Comme $W$ est localement noethérien, le morphisme
$w : \Spec(k) \to W$ est de présentation finie, voir
Morphismes, Lemme \ref{morphisms-lemma-closed-immersion-finite-presentation}.
Le composé
$$
\Spec(k) \xrightarrow{w} W \longrightarrow \mathcal{Z}
$$
est donc localement de présentation finie d'après
Morphismes d'espaces, Lemmes
\ref{spaces-morphisms-lemma-composition-finite-presentation} et
\ref{spaces-morphisms-lemma-smooth-locally-finite-presentation}.
Il est aussi plat et surjectif d'après le
Lemme \ref{lemma-flat-cover-by-field}.
Ainsi, (2) est vérifiée.
\end{proof}

\begin{lemma}
\label{lemma-monomorphism-into-point}
Soit $\mathcal{Z}' \to \mathcal{Z}$ un monomorphisme de champs algébriques.
Supposons qu'il existe un corps $k$ et un morphisme localement de présentation finie,
surjectif et plat $\Spec(k) \to \mathcal{Z}$. Alors soit $\mathcal{Z}'$
est vide, soit $\mathcal{Z}' \to \mathcal{Z}$ est une équivalence.
\end{lemma}

\begin{proof}
Nous pouvons supposer que $\mathcal{Z}'$ est non vide. Dans ce cas, le
produit fibré $T = \mathcal{Z}' \times_\mathcal{Z} \Spec(k)$
est non vide, voir
Lemme \ref{lemma-points-cartesian}.
Or $T$ est un espace algébrique et la projection $T \to \Spec(k)$
est un monomorphisme. Ainsi, $T = \Spec(k)$, voir
Morphismes d'espaces, Lemme
\ref{spaces-morphisms-lemma-monomorphism-toward-field}.
Nous concluons que $\Spec(k) \to \mathcal{Z}$ se factorise par
$\mathcal{Z}'$. Supposons que le morphisme $z : \Spec(k) \to \mathcal{Z}$
soit donné par l'objet $\xi$ au-dessus de $\Spec(k)$. Nous venons de voir que
$\xi$ est isomorphe à un objet $\xi'$ de $\mathcal{Z}'$ au-dessus de
$\Spec(k)$. Puisque $z$
est surjectif, plat et localement de présentation finie, on voit que
tout objet de $\mathcal{Z}$ au-dessus d'un schéma quelconque est localement pour fppf isomorphe
à un tiré en arrière de $\xi$, et donc aussi à un tiré en arrière de $\xi'$. Par descente des
objets pour les champs en groupoïdes, cela implique que
$\mathcal{Z}' \to \mathcal{Z}$ est essentiellement surjectif (et aussi
pleinement fidèle, voir
Lemme \ref{lemma-monomorphism}).
D'où le résultat.
\end{proof}

\begin{lemma}
\label{lemma-improve-unique-point}
Soit $\mathcal{Z}$ un champ algébrique. Supposons que $\mathcal{Z}$ satisfasse
aux conditions équivalentes du
Lemme \ref{lemma-unique-point}.
Il existe alors une unique sous-catégorie strictement pleine
$\mathcal{Z}' \subset \mathcal{Z}$ telle que
$\mathcal{Z}'$ soit un champ algébrique satisfaisant aux conditions équivalentes du
Lemme \ref{lemma-unique-point-better}.
Le morphisme d'inclusion $\mathcal{Z}' \to \mathcal{Z}$ est un monomorphisme
de champs algébriques.
\end{lemma}

\begin{proof}
La dernière assertion résulte immédiatement de la première et du
Lemme \ref{lemma-monomorphism}.
Choisissons un corps $k$ et un morphisme $\Spec(k) \to \mathcal{Z}$
qui soit surjectif, plat et localement de type fini.
Posons $U = \Spec(k)$ et $R = U \times_\mathcal{Z} U$.
Les projections $s, t : R \to U$ sont localement de type fini.
Comme $U$ est le spectre d'un corps, il s'ensuit que
$s, t$ sont plates et localement de présentation finie (d'après
Morphismes d'espaces, Lemme
\ref{spaces-morphisms-lemma-noetherian-finite-type-finite-presentation}).
On voit que $\mathcal{Z}' = [U/R]$ est un champ algébrique d'après
Critères de représentabilité,
Théorème \ref{criteria-theorem-flat-groupoid-gives-algebraic-stack}.
D'après
Champs algébriques, Lemme \ref{algebraic-lemma-map-space-into-stack},
nous obtenons un morphisme canonique
$$
f : \mathcal{Z}' \longrightarrow \mathcal{Z}
$$
qui est pleinement fidèle. Ce morphisme est donc représentable par des
espaces algébriques, voir
Champs algébriques, Lemme
\ref{algebraic-lemma-characterize-representable-by-algebraic-spaces},
et c'est un monomorphisme, voir
Lemme \ref{lemma-monomorphism}.
D'après
Critères de représentabilité,
Lemme \ref{criteria-lemma-flat-quotient-flat-presentation},
le morphisme $U \to \mathcal{Z}'$ est surjectif, plat et localement de
présentation finie. Ainsi, $\mathcal{Z}'$ est un champ algébrique qui satisfait
aux conditions équivalentes du
Lemme \ref{lemma-unique-point-better}.
D'après
Champs algébriques, Lemme \ref{algebraic-lemma-equivalent},
nous pouvons remplacer $\mathcal{Z}'$ par son image essentielle dans $\mathcal{Z}$.
Nous avons donc prouvé toutes les assertions du lemme, à l'exception de
l'unicité de $\mathcal{Z}' \subset \mathcal{Z}$. Supposons que
$\mathcal{Z}'' \subset \mathcal{Z}$ soit un second champ algébrique de ce type.
Alors les projections
$$
\mathcal{Z}'
\longleftarrow
\mathcal{Z}' \times_\mathcal{Z} \mathcal{Z}''
\longrightarrow
\mathcal{Z}''
$$
sont des monomorphismes. Le champ algébrique du milieu est non vide d'après le
Lemme \ref{lemma-points-cartesian}.
Les deux projections sont donc des isomorphismes d'après le
Lemme \ref{lemma-monomorphism-into-point},
d'où le résultat.
\end{proof}


\begin{example}
\label{example-distinct}
Voici un exemple où le morphisme construit dans le
Lemme \ref{lemma-improve-unique-point}
n'est pas un isomorphisme. Cet exemple montre qu'il est nécessaire d'imposer
aux gerbes résiduelles d'être localement noethériennes dans la
Définition \ref{definition-residual-gerbe}.
En fait, l'exemple est même un espace algébrique !
Soit $\text{Gal}(\overline{\mathbf{Q}}/\mathbf{Q})$ le groupe de
Galois absolu de $\mathbf{Q}$ muni de la topologie profinie.
Soit
$$
U =
\Spec(\overline{\mathbf{Q}})
\times_{\Spec(\mathbf{Q})}
\Spec(\overline{\mathbf{Q}}) =
\text{Gal}(\overline{\mathbf{Q}}/\mathbf{Q}) \times
\Spec(\overline{\mathbf{Q}})
$$
(nous omettons une explication précise du sens du dernier signe égal).
Notons $G$ le groupe de Galois absolu
$\text{Gal}(\overline{\mathbf{Q}}/\mathbf{Q})$ muni de la topologie
discrète, considéré comme un schéma en groupes constant sur
$\Spec(\overline{\mathbf{Q}})$, voir
Groupoïdes, Exemple \ref{groupoids-example-constant-group}.
Alors $G$ agit librement et transitivement sur $U$.
Soit $X = U/G$, voir
Espaces, Définition \ref{spaces-definition-quotient}.
Alors $X$ est un espace algébrique réduit non noethérien ayant exactement un point.
De plus, $X$ possède un point (localement) de type fini :
$$
x :
\Spec(\overline{\mathbf{Q}}) \longrightarrow
U \longrightarrow X
$$
En effet, tout point de $U$ est en fait fermé !
Comme $X$ est un espace algébrique sur $\overline{\mathbf{Q}}$,
il s'ensuit que $x$ est un monomorphisme. Ainsi, $x$ est le morphisme
construit dans le
Lemme \ref{lemma-improve-unique-point},
mais $x$ n'est pas un isomorphisme. En fait,
$\Spec(\overline{\mathbf{Q}}) \to X$ est la gerbe résiduelle de $X$ en $x$.
\end{example}

\noindent
Nous verrons plus loin que, sous des hypothèses modérées sur le champ algébrique
$\mathcal{X}$, les conditions équivalentes du lemme suivant sont satisfaites
pour tout point $x \in |\mathcal{X}|$ (voir Morphismes de champs, Section
\ref{stacks-morphisms-section-existence-residual-gerbes}).

\begin{lemma}
\label{lemma-residual-gerbe}
Soit $\mathcal{X}$ un champ algébrique. Soit $x \in |\mathcal{X}|$ un point.
Les assertions suivantes sont équivalentes :
\begin{enumerate}
\item il existe un champ algébrique $\mathcal{Z}$ et un monomorphisme
$\mathcal{Z} \to \mathcal{X}$ tels que $|\mathcal{Z}|$ soit réduit à un point
et que l'image de $|\mathcal{Z}|$ dans $|\mathcal{X}|$ soit $x$,
\item il existe un champ algébrique réduit $\mathcal{Z}$ et un monomorphisme
$\mathcal{Z} \to \mathcal{X}$ tels que $|\mathcal{Z}|$ soit réduit à un point
et que l'image de $|\mathcal{Z}|$ dans $|\mathcal{X}|$ soit $x$,
\item il existe un champ algébrique $\mathcal{Z}$, un monomorphisme
$f : \mathcal{Z} \to \mathcal{X}$ et un morphisme plat surjectif
$z : \Spec(k) \to \mathcal{Z}$, où $k$ est un corps tel que
$x = f(z)$.
\end{enumerate}
De plus, si ces conditions sont satisfaites, il existe une unique
sous-catégorie strictement pleine $\mathcal{Z}_x \subset \mathcal{X}$
telle que $\mathcal{Z}_x$ soit un champ algébrique réduit et localement
noethérien, et que $|\mathcal{Z}_x|$ soit réduit à un point envoyé sur $x$
par l'application $|\mathcal{Z}_x| \to |\mathcal{X}|$.
\end{lemma}

\begin{proof}
Si $\mathcal{Z} \to \mathcal{X}$ est comme en (1), alors
$\mathcal{Z}_{red} \to \mathcal{X}$ est comme en (2). (Voir la
Section \ref{section-reduced}
pour la notion de réduction d'un champ algébrique.)
Ainsi, (1) implique (2).
Il est immédiat que (2) implique (1).
L'équivalence de (2) et (3) résulte immédiatement du
Lemme \ref{lemma-unique-point}.

\medskip\noindent
À ce stade, nous avons établi l'équivalence de (1) -- (3).
Choisissons un monomorphisme $f : \mathcal{Z} \to \mathcal{X}$ comme en (2).
Remarquons qu'il en résulte que $f$ est pleinement fidèle, voir
Lemme \ref{lemma-monomorphism}.
Notons $\mathcal{Z}' \subset \mathcal{X}$ l'image essentielle du foncteur
$f$. Alors $f : \mathcal{Z} \to \mathcal{Z}'$ est une équivalence, et donc
$\mathcal{Z}'$ est un champ algébrique, voir
Champs algébriques, Lemme \ref{algebraic-lemma-equivalent}.
Appliquons le
Lemme \ref{lemma-improve-unique-point}
pour obtenir une sous-catégorie strictement pleine $\mathcal{Z}_x \subset \mathcal{Z}'$
comme dans l'énoncé du lemme.
Cela prouve toutes les assertions du lemme, sauf l'unicité.

\medskip\noindent
Pour prouver l'unicité, supposons que
$\mathcal{Z}_x \subset \mathcal{X}$
et
$\mathcal{Z}'_x \subset \mathcal{X}$
soient deux sous-catégories strictement pleines comme dans l'énoncé du lemme.
Alors les projections
$$
\mathcal{Z}'_x
\longleftarrow
\mathcal{Z}'_x \times_\mathcal{X} \mathcal{Z}_x
\longrightarrow
\mathcal{Z}_x
$$
sont des monomorphismes. Le champ algébrique du milieu est non vide d'après le
Lemme \ref{lemma-points-cartesian}.
Les deux projections sont donc des isomorphismes d'après le
Lemme \ref{lemma-monomorphism-into-point},
d'où le résultat.
\end{proof}

\noindent
Ces explications nous permettent maintenant de poser la définition suivante.

\begin{definition}
\label{definition-residual-gerbe}
Soit $\mathcal{X}$ un champ algébrique. Soit $x \in |\mathcal{X}|$.
\begin{enumerate}
\item On dit que la {\it gerbe résiduelle de $\mathcal{X}$ en $x$ existe}
si les conditions équivalentes (1), (2) et (3) du
Lemme \ref{lemma-residual-gerbe}
sont satisfaites.
\item Si la gerbe résiduelle de $\mathcal{X}$ en $x$ existe, la
{\it gerbe résiduelle de $\mathcal{X}$ en $x$}\footnote{Cela entre en conflit avec
\cite{LM-B} dans l'esprit, mais non en fait. Plus précisément, au chapitre 11, ils associent
à tout point de tout champ algébrique quasi-séparé une gerbe (pas nécessairement
algébrique), qu'ils appellent la gerbe résiduelle. Nous verrons dans
Morphismes de champs, Lemme
\ref{stacks-morphisms-lemma-every-point-residual-gerbe}
que, sur un champ algébrique quasi-séparé, tout point
possède une gerbe résiduelle en notre sens, qui est alors équivalente à la leur. Pour
plus de renseignements sur ce sujet, voir
\cite[Annexe B]{rydh_etale_devissage}.}
est la sous-catégorie strictement pleine
$\mathcal{Z}_x \subset \mathcal{X}$ construite au
Lemme \ref{lemma-residual-gerbe}.
\end{enumerate}
\end{definition}

\noindent
En particulier, nous savons que $\mathcal{Z}_x$ (si elle existe) est un
champ algébrique localement noethérien et réduit, et qu'il existe un
corps et un morphisme surjectif, plat et localement de présentation finie
$$
\Spec(k) \longrightarrow \mathcal{Z}_x.
$$
Nous verrons dans
Morphismes de champs, Lemme
\ref{stacks-morphisms-lemma-noetherian-singleton-stack-gerbe}
que $\mathcal{Z}_x$ est une gerbe. L'existence des gerbes résiduelles est
étudiée dans Morphismes de champs, Section
\ref{stacks-morphisms-section-existence-residual-gerbes}.

\begin{example}
\label{example-residual-gerbe}
Soit $X$ un schéma et soit $x \in X$ un point. Alors le monomorphisme
$x \to X$ est la gerbe résiduelle de $X$ en $x$, où, comme d'habitude, nous identifions
$x$ au schéma $x = \Spec(\kappa(x))$. Si $X$ est un espace algébrique
et $x \in |X|$, alors la gerbe résiduelle en $x$ (appelée
l'espace résiduel) existe toujours, voir
Espaces décents, Section \ref{decent-spaces-section-singleton}.
\end{example}

\noindent
La gerbe résiduelle, si elle existe, est un champ algébrique régulier
d'après le lemme suivant.

\begin{lemma}
\label{lemma-residual-gerbe-regular}
Un champ algébrique réduit et localement noethérien $\mathcal{Z}$ tel que
$|\mathcal{Z}|$ soit réduit à un point est régulier.
\end{lemma}

\begin{proof}
Soit $W \to \mathcal{Z}$ un morphisme lisse surjectif, où $W$ est un schéma.
Soit $k$ un corps et soit $\Spec(k) \to \mathcal{Z}$ un morphisme surjectif,
plat et localement de présentation finie (voir
Lemme \ref{lemma-unique-point-better}).
L'espace algébrique $T = W \times_\mathcal{Z} \Spec(k)$ est
lisse sur $k$, donc en particulier régulier, voir
Espaces sur les corps, Lemme \ref{spaces-over-fields-lemma-smooth-regular}.
Comme $T \to W$ est localement de présentation finie, plat et surjectif, il
s'ensuit que $W$ est régulier, voir
Descente sur les espaces, Lemme \ref{spaces-descent-lemma-descend-regular}.
Par définition, cela signifie que $\mathcal{Z}$ est régulier.
\end{proof}

\begin{lemma}
\label{lemma-residual-gerbe-points}
Soit $\mathcal{X}$ un champ algébrique. Soit $x \in |\mathcal{X}|$.
Supposons que la gerbe résiduelle $\mathcal{Z}_x$ de $\mathcal{X}$ existe.
Soit $f : \Spec(K) \to \mathcal{X}$ un morphisme, où $K$ est un corps,
dans la classe d'équivalence de $x$. Alors $f$ se factorise par le morphisme
d'inclusion $\mathcal{Z}_x \to \mathcal{X}$.
\end{lemma}

\begin{proof}
Choisissons un corps $k$ et un morphisme surjectif, plat et localement de présentation finie
$\Spec(k) \to \mathcal{Z}_x$. Posons
$T = \Spec(K) \times_\mathcal{X} \mathcal{Z}_x$. D'après le
Lemme \ref{lemma-points-cartesian},
on voit que $T$ est non vide. Comme $\mathcal{Z}_x \to \mathcal{X}$
est un monomorphisme, on voit que $T \to \Spec(K)$ est un monomorphisme.
Ainsi, d'après
Morphismes d'espaces, Lemme
\ref{spaces-morphisms-lemma-monomorphism-toward-field},
on voit que $T = \Spec(K)$, ce qui prouve le lemme.
\end{proof}

\begin{lemma}
\label{lemma-residual-gerbe-unique}
Soit $\mathcal{X}$ un champ algébrique. Soit $x \in |\mathcal{X}|$.
Soit $\mathcal{Z}$ un champ algébrique satisfaisant aux conditions équivalentes du
Lemme \ref{lemma-unique-point-better},
et soit $\mathcal{Z} \to \mathcal{X}$ un monomorphisme tel que l'image
de $|\mathcal{Z}| \to |\mathcal{X}|$ soit $x$. Alors la gerbe résiduelle
$\mathcal{Z}_x$ de $\mathcal{X}$ en $x$ existe, et
$\mathcal{Z} \to \mathcal{X}$ se factorise sous la forme
$\mathcal{Z} \to \mathcal{Z}_x \to \mathcal{X}$, où la première flèche
est une équivalence.
\end{lemma}

\begin{proof}
Soit $\mathcal{Z}_x \subset \mathcal{X}$ la sous-catégorie pleine correspondant
à l'image essentielle du foncteur $\mathcal{Z} \to \mathcal{X}$.
Alors $\mathcal{Z} \to \mathcal{Z}_x$ est une équivalence, donc
$\mathcal{Z}_x$ est un champ algébrique, voir
Champs algébriques, Lemme \ref{algebraic-lemma-equivalent}.
Comme $\mathcal{Z}_x$ hérite de toutes les propriétés de $\mathcal{Z}$ par
cette équivalence, il résulte clairement de l'unicité du
Lemme \ref{lemma-residual-gerbe}
que $\mathcal{Z}_x$ est la gerbe résiduelle de $\mathcal{X}$ en $x$.
\end{proof}

\begin{lemma}
\label{lemma-residual-gerbe-functorial}
Soit $f : \mathcal{X} \to \mathcal{Y}$ un morphisme de champs algébriques.
Soit $x \in |\mathcal{X}|$, d'image $y \in |\mathcal{Y}|$.
Si les gerbes résiduelles $\mathcal{Z}_x \subset \mathcal{X}$
et $\mathcal{Z}_y \subset \mathcal{Y}$ de $x$ et de $y$ existent,
alors $f$ induit un diagramme commutatif
$$
\xymatrix{
\mathcal{X} \ar[d]_f & \mathcal{Z}_x \ar[l] \ar[d] \\
\mathcal{Y} & \mathcal{Z}_y \ar[l]
}
$$
\end{lemma}

\begin{proof}
Choisissons un corps $k$ et un morphisme surjectif, plat et localement de présentation finie
$\Spec(k) \to \mathcal{Z}_x$. Le morphisme
$\Spec(k) \to \mathcal{Y}$ se factorise par $\mathcal{Z}_y$ d'après le
Lemme \ref{lemma-residual-gerbe-points}.
Ainsi, $\mathcal{Z}_x \times_\mathcal{Y} \mathcal{Z}_y$
est un sous-champ non vide de $\mathcal{Z}_x$,
donc égal à $\mathcal{Z}_x$ d'après le Lemme \ref{lemma-monomorphism-into-point}.
\end{proof}

\begin{lemma}
\label{lemma-residual-gerbe-isomorphic}
Soit $f : \mathcal{X} \to \mathcal{Y}$ un morphisme de champs algébriques.
Soit $x \in |\mathcal{X}|$ dont l'image est $y \in |\mathcal{Y}|$.
Supposons que les gerbes résiduelles $\mathcal{Z}_x \subset \mathcal{X}$
et $\mathcal{Z}_y \subset \mathcal{Y}$ de $x$ et de $y$ existent
et qu'il existe un morphisme $\Spec(k) \to \mathcal{X}$
dans la classe d'équivalence de $x$ tel que
$$
\Spec(k) \times_\mathcal{X} \Spec(k)
\longrightarrow
\Spec(k) \times_\mathcal{Y} \Spec(k)
$$
soit un isomorphisme. Alors $\mathcal{Z}_x \to \mathcal{Z}_y$
est un isomorphisme.
\end{lemma}

\begin{proof}
Soit $k'/k$ une extension de corps. Alors
$$
\Spec(k') \times_\mathcal{X} \Spec(k')
\longrightarrow
\Spec(k') \times_\mathcal{Y} \Spec(k')
$$
est le changement de base du morphisme du lemme par le
morphisme fidèlement plat $\Spec(k' \otimes k') \to \Spec(k \otimes k)$.
Ainsi, la propriété décrite dans le lemme ne dépend pas du
choix du morphisme $\Spec(k) \to \mathcal{X}$ dans la
classe d'équivalence de $x$. Ainsi, nous pouvons supposer que
$\Spec(k) \to \mathcal{Z}_x$ est surjectif, plat et
localement de présentation finie. Dans cette situation, nous avons
$$
\mathcal{Z}_x = [\Spec(k)/R]
$$
où $R = \Spec(k) \times_\mathcal{X} \Spec(k)$. Voir la
démonstration du Lemme \ref{lemma-improve-unique-point}.
Comme on a aussi $R = \Spec(k) \times_\mathcal{Y} \Spec(k)$,
nous concluons que le morphisme $\mathcal{Z}_x \to \mathcal{Z}_y$
du Lemme \ref{lemma-residual-gerbe-functorial}
est pleinement fidèle d'après le
Lemme sur les champs algébriques \ref{algebraic-lemma-map-space-into-stack}.
On conclut, par exemple, à l'aide du Lemme \ref{lemma-residual-gerbe-unique}.
\end{proof}








\section{Dimension d'un champ}
\label{section-dimension}

\noindent
Nous pouvons définir la dimension d'un champ algébrique $\mathcal{X}$
en un point $x$ à l'aide de la notion de dimension d'un espace algébrique
en un point (Propriétés des espaces, Définition
\ref{spaces-properties-definition-dimension-at-point}).
Dans le lemme suivant, le résultat peut être $\infty$, soit parce que
$\mathcal{X}$ n'est pas quasi-compact, soit parce que nous rencontrons le
phénomène décrit dans Exemples, Section
\ref{examples-section-Noetherian-infinite-dimension}.

\begin{lemma}
\label{lemma-dimension-at-point-well-defined}
Soit $\mathcal{X}$ un champ algébrique localement noethérien sur un schéma $S$.
Soit $x \in |\mathcal{X}|$ un point de $\mathcal{X}$.
Soit $[U/R] \to \mathcal{X}$ une présentation
(Champs algébriques, Définition \ref{algebraic-definition-presentation})
où $U$ est un schéma. Soit $u \in U$ un point dont l'image est $x$.
Soit $e : U \to R$ le morphisme « identité » et soit $s : R \to U$ le
morphisme « source », qui est un morphisme lisse d'espaces algébriques. Soit $R_u$
la fibre de $s : R \to U$ au-dessus de $u$. L'élément
$$
\dim_x(\mathcal{X}) = \dim_u(U) - \dim_{e(u)}(R_u) \in \mathbf{Z} \cup \infty
$$
est indépendant du choix de la présentation et du point $u$ au-dessus de $x$.
\end{lemma}

\begin{proof}
Puisque $R \to U$ est lisse, le schéma $R_u$ est lisse sur $\kappa(u)$
et est donc de dimension finie. D'autre part, le schéma $U$
est localement noethérien, mais cela ne garantit pas que
$\dim_u(U)$ soit finie. Ainsi, la différence est un élément de
$\mathbf{Z} \cup \{\infty\}$.

\medskip\noindent
Soient $[U'/R'] \to \mathcal{X}$ et $u' \in U'$ une seconde présentation
où $U'$ est un schéma et où l'image de $u'$ est $x$. Considérons l'espace
algébrique $P = U \times_\mathcal{X} U'$. D'après le
Lemme \ref{lemma-points-cartesian}, il existe $p \in |P|$ dont les images sont
$u$ et $u'$. Comme $P \to U$ et $P \to U'$ sont lisses, on voit que
$\dim_p(P) = \dim_u(U) + \dim_p(P_u)$ et
$\dim_p(P) = \dim_{u'}(U') + \dim_p(P_{u'})$; voir le
Lemme de Morphismes d'espaces
\ref{spaces-morphisms-lemma-smoothness-dimension-spaces}.
Notons que
$$
R'_{u'} = \Spec(\kappa(u')) \times_\mathcal{X} U'
\quad\text{and}\quad
P_u = \Spec(\kappa(u)) \times_\mathcal{X} U'
$$
Représentons $p \in |P|$ par un morphisme $\Spec(\Omega) \to P$.
Comme $p$ s'envoie à la fois sur $u$ et sur $u'$, il induit un $2$-morphisme
entre les composées
$\Spec(\Omega) \to \Spec(\kappa(u')) \to \mathcal{X}$ et
$\Spec(\Omega) \to \Spec(\kappa(u)) \to \mathcal{X}$,
qui définit à son tour un isomorphisme
$$
\Spec(\Omega) \times_{\Spec(\kappa(u'))} R'_{u'}
\cong
\Spec(\Omega) \times_{\Spec(\kappa(u))} P_u
$$
d'espaces algébriques sur $\Spec(\Omega)$ envoyant le point $\Omega$-rationnel
$(1, e'(u'))$ sur $(1, p)$ (certains détails sont omis). On en déduit que
$$
\dim_{e'(u')}(R'_{u'}) = \dim_p(P_u)
$$
d'après le Lemme de Morphismes d'espaces
\ref{spaces-morphisms-lemma-dimension-fibre-after-base-change}.
Par symétrie, on a
$\dim_{e(u)}(R_u) = \dim_p(P_{u'})$.
En réunissant toutes ces égalités, on obtient l'indépendance des choix.
\end{proof}

\noindent
Nous pouvons utiliser le lemme précédent pour poser la définition suivante.

\begin{definition}
\label{definition-dimension-at-point}
Soit $\mathcal{X}$ un champ algébrique localement noethérien sur un schéma $S$.
Soit $x \in |\mathcal{X}|$ un point de $\mathcal{X}$.
Soit $[U/R] \to \mathcal{X}$ une présentation
(Champs algébriques, Définition \ref{algebraic-definition-presentation})
où $U$ est un schéma,
et soit $u \in U$ un point dont l'image est $x$.
Nous définissons la {\it dimension de $\mathcal{X}$ en $x$} comme
l'élément $\dim_x(\mathcal{X}) \in \mathbf{Z} \cup \infty$
tel que 
$$
\dim_x(\mathcal{X}) = \dim_u(U)-\dim_{e(u)}(R_u).
$$
Les notations sont celles du Lemme \ref{lemma-dimension-at-point-well-defined}.
\end{definition}

\noindent
La dimension d'un champ en un point coïncide avec la notion usuelle
lorsque $\mathcal{X}$ est un schéma (Topologie, Définition
\ref{topology-definition-Krull}),
ou, plus généralement, lorsque $\mathcal{X}$ est un espace algébrique localement noethérien
(Propriétés des espaces, Définition
\ref{spaces-properties-definition-dimension-at-point}).

\begin{definition}
\label{definition-dimension}
Soit $S$ un schéma. Soit $\mathcal{X}$
un champ algébrique localement noethérien sur $S$.
La {\it dimension} $\dim(\mathcal{X})$ de $\mathcal{X}$ est définie par
$$
\dim(\mathcal{X}) = \sup\nolimits_{x \in |\mathcal{X}|} \dim_x(\mathcal{X})
$$
\end{definition}

\noindent
Cette définition de la dimension coïncide avec la notion usuelle
si $\mathcal{X}$ est un schéma
(Propriétés, Lemme \ref{properties-lemma-dimension})
ou un espace algébrique (Propriétés des espaces, Définition
\ref{spaces-properties-definition-dimension}).

\begin{remark}
\label{remark-dimension-empty-stack}
Si $\mathcal{X}$ est un champ non vide de type fini sur un corps,
alors $\dim(\mathcal{X})$ est un entier. Pour un champ algébrique
localement noethérien quelconque $\mathcal{X}$,
$\dim(\mathcal{X})$ appartient à $Z\cup \{\pm \infty\}$,
et $\dim(\mathcal{X}) = -\infty$ si et seulement si $\mathcal{X}$
est vide.
\end{remark}

\begin{example}
\label{example-dimension-quotient-stack}
Soit $X$ un schéma de type fini sur un corps $k$, et soit $G$
un schéma en groupes de type fini sur $k$ qui opère sur $X$.
Alors la dimension du champ quotient $[X/G]$ est égale à
$\dim(X)-\dim(G)$. En particulier, la dimension du champ classifiant
$BG=[\Spec(k)/G]$ est $-\dim(G)$. Ainsi, la dimension d'un champ algébrique
peut être un entier négatif, contrairement à ce qui se produit pour les schémas
ou les espaces algébriques.
\end{example}






\section{Irréductibilité locale}
\label{section-irreducible-local-ring}

\noindent
Nous avons défini le nombre de branches géométriques d'un schéma en un point
dans Propriétés, Section \ref{properties-section-unibranch},
et celui d'un espace algébrique en un point dans Propriétés des espaces, Section
\ref{spaces-properties-section-irreducible-local-ring}.
Soit $n \in \mathbf{N}$. Pour un anneau local $A$, posons
$$
P_n(A) = \text{the number of geometric branches of }A\text{ is }n
$$
Pour un homomorphisme lisse d'anneaux $A \to B$ et un idéal premier $\mathfrak q$
de $B$ situé au-dessus de $\mathfrak p$ de $A$, on a
$$
P_n(A_\mathfrak p) \Leftrightarrow P_n(B_\mathfrak q)
$$
d'après le Lemme de Compléments d'algèbre
\ref{more-algebra-lemma-invariance-number-branches-smooth}.
Comme dans Propriétés des espaces, Remarque
\ref{spaces-properties-remark-list-properties-local-ring-local-etale-topology},
nous pouvons utiliser $P_n$ pour définir une propriété locale pour la topologie \'etale $\mathcal{P}_n$
des germes $(U, u)$ de schémas en posant
$\mathcal{P}_n(U, u) = P_n(\mathcal{O}_{U, u})$.
La propriété correspondante $\mathcal{P}_n$
d'un espace algébrique $X$ en un point $x$
(voir Propriétés des espaces, Définition
\ref{spaces-properties-definition-property-at-point})
est précisément la propriété
``le nombre de branches géométriques de $X$ en $x$ est $n$''; voir
Propriétés des espaces, Définition
\ref{spaces-properties-definition-number-of-geometric-branches}.
De plus, la propriété $\mathcal{P}_n$ est locale pour la topologie lisse; voir
Descente, Définition \ref{descent-definition-local-at-point}.
Cela résulte soit de l'équivalence affichée ci-dessus,
soit du Lemme de Compléments sur les morphismes
\ref{more-morphisms-lemma-number-of-branches-and-smooth}.
Ainsi, la Définition \ref{definition-property-at-point}
s'applique et nous obtenons une notion pour les champs algébriques en un point.

\begin{definition}
\label{definition-number-of-geometric-branches}
Soit $\mathcal{X}$ un champ algébrique. Soit $x \in |\mathcal{X}|$.
\begin{enumerate}
\item Le {\it nombre de branches géométriques de $\mathcal{X}$ en $x$}
vaut soit $n \in \mathbf{N}$ si les conditions équivalentes du
Lemme \ref{lemma-local-source-target-at-point} sont satisfaites pour
$\mathcal{P}_n$ définie ci-dessus, soit $\infty$.
\item Nous disons que $\mathcal{X}$ est {\it géométriquement unibranche en $x$}
si le nombre de branches géométriques de $\mathcal{X}$ en $x$ vaut $1$.
\end{enumerate}
\end{definition}






\section{Conditions de finitude et points}
\label{section-points-conditions}

\noindent
Cette section est l'analogue de
Espaces décents, Section \ref{decent-spaces-section-points-monomorphisms}
pour les points des champs algébriques.

\begin{lemma}
\label{lemma-UR-quasi-compact-above-x}
Soit $\mathcal{X}$ un champ algébrique. Soit $x \in |\mathcal{X}|$
un point. Les conditions suivantes sont équivalentes :
\begin{enumerate}
\item un morphisme $\Spec(k) \to \mathcal{X}$ dans la classe d'équivalence
de $x$ est quasi-compact, et
\item tout morphisme $\Spec(k) \to \mathcal{X}$ dans la classe d'équivalence
de $x$ est quasi-compact.
\end{enumerate}
\end{lemma}

\begin{proof}
Soit $\Spec(k) \to \mathcal{X}$ dans la classe d'équivalence
de $x$. Soit $k'/k$ une extension de corps.
Nous devons alors montrer que $\Spec(k) \to \mathcal{X}$ est
quasi-compact si et seulement si $\Spec(k') \to \mathcal{X}$
est quasi-compact. Cela résulte du
Lemme de Morphismes d'espaces
\ref{spaces-morphisms-lemma-surjection-from-quasi-compact}
et du principe donné par le Lemme de Champs algébriques
\ref{algebraic-lemma-representable-transformations-property-implication}.
\end{proof}

\noindent
On dit parfois qu'un point $x \in |\mathcal{X}|$ satisfaisant
aux conditions équivalentes du Lemme \ref{lemma-UR-quasi-compact-above-x}
est un ``{\it point quasi-compact}''.









\begin{multicols}{2}[\section{Autres chapitres}]
\noindent
Préliminaires
\begin{enumerate}
\item \hyperref[introduction-section-phantom]{Introduction}
\item \hyperref[conventions-section-phantom]{Conventions}
\item \hyperref[sets-section-phantom]{Théorie des ensembles}
\item \hyperref[categories-section-phantom]{Catégories}
\item \hyperref[topology-section-phantom]{Topologie}
\item \hyperref[sheaves-section-phantom]{Faisceaux sur les espaces}
\item \hyperref[sites-section-phantom]{Sites et faisceaux}
\item \hyperref[stacks-section-phantom]{Champs}
\item \hyperref[fields-section-phantom]{Corps}
\item \hyperref[algebra-section-phantom]{Algèbre commutative}
\item \hyperref[brauer-section-phantom]{Groupes de Brauer}
\item \hyperref[homology-section-phantom]{Algèbre homologique}
\item \hyperref[derived-section-phantom]{Catégories dérivées}
\item \hyperref[simplicial-section-phantom]{Méthodes simpliciales}
\item \hyperref[more-algebra-section-phantom]{Compléments d'algèbre}
\item \hyperref[smoothing-section-phantom]{Lissage des morphismes d'anneaux}
\item \hyperref[modules-section-phantom]{Faisceaux de modules}
\item \hyperref[sites-modules-section-phantom]{Modules sur les sites}
\item \hyperref[injectives-section-phantom]{Objets injectifs}
\item \hyperref[cohomology-section-phantom]{Cohomologie des faisceaux}
\item \hyperref[sites-cohomology-section-phantom]{Cohomologie sur les sites}
\item \hyperref[dga-section-phantom]{Algèbre différentielle graduée}
\item \hyperref[dpa-section-phantom]{Algèbre à puissances divisées}
\item \hyperref[sdga-section-phantom]{Faisceaux différentiels gradués}
\item \hyperref[hypercovering-section-phantom]{Hyperrecouvrements}
\end{enumerate}
Schémas
\begin{enumerate}
\setcounter{enumi}{25}
\item \hyperref[schemes-section-phantom]{Schémas}
\item \hyperref[constructions-section-phantom]{Constructions de schémas}
\item \hyperref[properties-section-phantom]{Propriétés des schémas}
\item \hyperref[morphisms-section-phantom]{Morphismes de schémas}
\item \hyperref[coherent-section-phantom]{Cohomologie des schémas}
\item \hyperref[divisors-section-phantom]{Diviseurs}
\item \hyperref[limits-section-phantom]{Limites de schémas}
\item \hyperref[varieties-section-phantom]{Variétés}
\item \hyperref[topologies-section-phantom]{Topologies sur les schémas}
\item \hyperref[descent-section-phantom]{Descente}
\item \hyperref[perfect-section-phantom]{Catégories dérivées de schémas}
\item \hyperref[more-morphisms-section-phantom]{Compléments sur les morphismes}
\item \hyperref[flat-section-phantom]{Compléments sur la platitude}
\item \hyperref[groupoids-section-phantom]{Schémas en groupoïdes}
\item \hyperref[more-groupoids-section-phantom]{Compléments sur les schémas en groupoïdes}
\item \hyperref[etale-section-phantom]{Morphismes étales de schémas}
\end{enumerate}
Thèmes de théorie des schémas
\begin{enumerate}
\setcounter{enumi}{41}
\item \hyperref[chow-section-phantom]{Homologie de Chow}
\item \hyperref[intersection-section-phantom]{Théorie de l'intersection}
\item \hyperref[pic-section-phantom]{Schémas de Picard des courbes}
\item \hyperref[weil-section-phantom]{Théories cohomologiques de Weil}
\item \hyperref[adequate-section-phantom]{Modules adéquats}
\item \hyperref[dualizing-section-phantom]{Complexes dualisants}
\item \hyperref[duality-section-phantom]{Dualité pour les schémas}
\item \hyperref[discriminant-section-phantom]{Discriminants et différentes}
\item \hyperref[derham-section-phantom]{Cohomologie de de Rham}
\item \hyperref[local-cohomology-section-phantom]{Cohomologie locale}
\item \hyperref[algebraization-section-phantom]{Géométrie algébrique et formelle}
\item \hyperref[curves-section-phantom]{Courbes algébriques}
\item \hyperref[resolve-section-phantom]{Résolution des surfaces}
\item \hyperref[models-section-phantom]{Réduction semi-stable}
\item \hyperref[functors-section-phantom]{Foncteurs et morphismes}
\item \hyperref[equiv-section-phantom]{Catégories dérivées de variétés}
\item \hyperref[pione-section-phantom]{Groupes fondamentaux de schémas}
\item \hyperref[etale-cohomology-section-phantom]{Cohomologie étale}
\item \hyperref[crystalline-section-phantom]{Cohomologie cristalline}
\item \hyperref[proetale-section-phantom]{Cohomologie pro-étale}
\item \hyperref[relative-cycles-section-phantom]{Cycles relatifs}
\item \hyperref[more-etale-section-phantom]{Compléments de cohomologie étale}
\item \hyperref[trace-section-phantom]{La formule des traces}
\end{enumerate}
Espaces algébriques
\begin{enumerate}
\setcounter{enumi}{64}
\item \hyperref[spaces-section-phantom]{Espaces algébriques}
\item \hyperref[spaces-properties-section-phantom]{Propriétés des espaces algébriques}
\item \hyperref[spaces-morphisms-section-phantom]{Morphismes d'espaces algébriques}
\item \hyperref[decent-spaces-section-phantom]{Espaces algébriques décents}
\item \hyperref[spaces-cohomology-section-phantom]{Cohomologie des espaces algébriques}
\item \hyperref[spaces-limits-section-phantom]{Limites d'espaces algébriques}
\item \hyperref[spaces-divisors-section-phantom]{Diviseurs sur les espaces algébriques}
\item \hyperref[spaces-over-fields-section-phantom]{Espaces algébriques sur des corps}
\item \hyperref[spaces-topologies-section-phantom]{Topologies sur les espaces algébriques}
\item \hyperref[spaces-descent-section-phantom]{Descente et espaces algébriques}
\item \hyperref[spaces-perfect-section-phantom]{Catégories dérivées d'espaces}
\item \hyperref[spaces-more-morphisms-section-phantom]{Compléments sur les morphismes d'espaces}
\item \hyperref[spaces-flat-section-phantom]{Platitude sur les espaces algébriques}
\item \hyperref[spaces-groupoids-section-phantom]{Groupoïdes dans les espaces algébriques}
\item \hyperref[spaces-more-groupoids-section-phantom]{Compléments sur les groupoïdes dans les espaces}
\item \hyperref[bootstrap-section-phantom]{Amorçage}
\item \hyperref[spaces-pushouts-section-phantom]{Sommes amalgamées d'espaces algébriques}
\end{enumerate}
Thèmes de géométrie
\begin{enumerate}
\setcounter{enumi}{81}
\item \hyperref[spaces-chow-section-phantom]{Groupes de Chow des espaces}
\item \hyperref[groupoids-quotients-section-phantom]{Quotients de groupoïdes}
\item \hyperref[spaces-more-cohomology-section-phantom]{Compléments sur la cohomologie des espaces}
\item \hyperref[spaces-simplicial-section-phantom]{Espaces simpliciaux}
\item \hyperref[spaces-duality-section-phantom]{Dualité pour les espaces}
\item \hyperref[formal-spaces-section-phantom]{Espaces algébriques formels}
\item \hyperref[restricted-section-phantom]{Algébrisation des espaces formels}
\item \hyperref[spaces-resolve-section-phantom]{Résolution des surfaces revisitée}
\end{enumerate}
Théorie des déformations
\begin{enumerate}
\setcounter{enumi}{89}
\item \hyperref[formal-defos-section-phantom]{Théorie formelle des déformations}
\item \hyperref[defos-section-phantom]{Théorie des déformations}
\item \hyperref[cotangent-section-phantom]{Le complexe cotangent}
\item \hyperref[examples-defos-section-phantom]{Problèmes de déformation}
\end{enumerate}
Champs algébriques
\begin{enumerate}
\setcounter{enumi}{93}
\item \hyperref[algebraic-section-phantom]{Champs algébriques}
\item \hyperref[examples-stacks-section-phantom]{Exemples de champs}
\item \hyperref[stacks-sheaves-section-phantom]{Faisceaux sur les champs algébriques}
\item \hyperref[criteria-section-phantom]{Critères de représentabilité}
\item \hyperref[artin-section-phantom]{Axiomes d'Artin}
\item \hyperref[quot-section-phantom]{Espaces de Quot et de Hilbert}
\item \hyperref[stacks-properties-section-phantom]{Propriétés des champs algébriques}
\item \hyperref[stacks-morphisms-section-phantom]{Morphismes de champs algébriques}
\item \hyperref[stacks-limits-section-phantom]{Limites de champs algébriques}
\item \hyperref[stacks-cohomology-section-phantom]{Cohomologie des champs algébriques}
\item \hyperref[stacks-perfect-section-phantom]{Catégories dérivées de champs}
\item \hyperref[stacks-introduction-section-phantom]{Introduction aux champs algébriques}
\item \hyperref[stacks-more-morphisms-section-phantom]{Compléments sur les morphismes de champs}
\item \hyperref[stacks-geometry-section-phantom]{La géométrie des champs}
\end{enumerate}
Thèmes de théorie des espaces de modules
\begin{enumerate}
\setcounter{enumi}{107}
\item \hyperref[moduli-section-phantom]{Champs de modules}
\item \hyperref[moduli-curves-section-phantom]{Espaces de modules de courbes}
\end{enumerate}
Divers
\begin{enumerate}
\setcounter{enumi}{109}
\item \hyperref[examples-section-phantom]{Exemples}
\item \hyperref[exercises-section-phantom]{Exercices}
\item \hyperref[guide-section-phantom]{Guide bibliographique}
\item \hyperref[desirables-section-phantom]{Desiderata}
\item \hyperref[coding-section-phantom]{Style de programmation}
\item \hyperref[obsolete-section-phantom]{Obsolète}
\item \hyperref[fdl-section-phantom]{Licence de documentation libre GNU}
\item \hyperref[index-section-phantom]{Index généré automatiquement}
\end{enumerate}
\end{multicols}


\bibliography{my}
\bibliographystyle{amsalpha}

\end{document}
